% ══════════════════════════════════════════════════════════════════════════════
% PREÁMBULO DEL LABORATORIO DE R PARA CIENCIA POLÍTICA — ITAM
% Emiliano Miranda González · Otoño 2026
%
% Estética: beamer académico de ciencia política internacional, en la tradición
% del deck de seminario (banda de color en la portada, láminas de sección a
% color completo con el índice atenuado, cuerpo blanco con el título del frame
% en el color de acento y sin regla, revelación progresiva de viñetas,
% pie discreto con sección y folio). Verde institucional #006847.
%
% Se usa así, desde la carpeta de una sesión:
%   % ══════════════════════════════════════════════════════════════════════════════
% PREÁMBULO DEL LABORATORIO DE R PARA CIENCIA POLÍTICA — ITAM
% Emiliano Miranda González · Otoño 2026
%
% Estética: beamer académico de ciencia política internacional, en la tradición
% del deck de seminario (banda de color en la portada, láminas de sección a
% color completo con el índice atenuado, cuerpo blanco con el título del frame
% en el color de acento y sin regla, revelación progresiva de viñetas,
% pie discreto con sección y folio). Verde institucional #006847.
%
% Se usa así, desde la carpeta de una sesión:
%   % ══════════════════════════════════════════════════════════════════════════════
% PREÁMBULO DEL LABORATORIO DE R PARA CIENCIA POLÍTICA — ITAM
% Emiliano Miranda González · Otoño 2026
%
% Estética: beamer académico de ciencia política internacional, en la tradición
% del deck de seminario (banda de color en la portada, láminas de sección a
% color completo con el índice atenuado, cuerpo blanco con el título del frame
% en el color de acento y sin regla, revelación progresiva de viñetas,
% pie discreto con sección y folio). Verde institucional #006847.
%
% Se usa así, desde la carpeta de una sesión:
%   % ══════════════════════════════════════════════════════════════════════════════
% PREÁMBULO DEL LABORATORIO DE R PARA CIENCIA POLÍTICA — ITAM
% Emiliano Miranda González · Otoño 2026
%
% Estética: beamer académico de ciencia política internacional, en la tradición
% del deck de seminario (banda de color en la portada, láminas de sección a
% color completo con el índice atenuado, cuerpo blanco con el título del frame
% en el color de acento y sin regla, revelación progresiva de viñetas,
% pie discreto con sección y folio). Verde institucional #006847.
%
% Se usa así, desde la carpeta de una sesión:
%   \input{../../estilo/preambulo-lab.tex}
%
% Se compila con pdflatex, DOS VECES (los nodos de tikz con `remember picture`
% necesitan la segunda pasada para colocarse bien).
% ══════════════════════════════════════════════════════════════════════════════

\documentclass[
  aspectratio=43,
  t,
  10pt
]{beamer}

% ── TIPOGRAFÍA ────────────────────────────────────────────────────────────────
% El cuerpo del deck sigue siendo sans en su totalidad. TeX Gyre Heros es un
% clon de Helvetica en Type 1 con cobertura T1 completa —incluidas cursivas y
% versalitas reales—, presente en cualquier distribución de TeX Live. lmodern
% va ANTES para quedarse con la monoespaciada: sin él, los comentarios en
% cursiva de los bloques de código salen como fuente de mapa de bits y se ven
% borrosos al proyectar.
%
% Si en la máquina de quien compila están instaladas Fira Sans o Lato —que son
% las humanistas que suelen usar los decks de seminario de ciencia política—,
% basta sustituir la línea de tgheros por \usepackage[sfdefault]{FiraSans} o
% \usepackage[default]{lato}. El resto del preámbulo no cambia.
\usepackage[T1]{fontenc}
\usepackage[utf8]{inputenc}
\usepackage{lmodern}
\usepackage{tgheros}
\renewcommand{\familydefault}{\sfdefault}
\usepackage{microtype}

% ── SERIFA PARA TÍTULOS (mezcla del deck de seminario) ─────────────────────────
% EMG pidió la mezcla de los decks de seminario de ciencia política
% internacional (ref.: deck de Gary King, IQSS/Harvard): serifa para lo que
% NOMBRA —portada y título de cada lámina— y sans para lo que INFORMA —cuerpo,
% viñetas, índice de secciones y pie—. Ver \serifalab más abajo y su uso en
% \setbeamerfont{title/subtitle/frametitle/framesubtitle}.
%
% Elección de la familia, verificada y no adivinada. De las candidatas del
% encargo, kpsewhich confirmó instaladas: tgpagella, tgtermes, tgschola,
% tgbonum, mathpazo, charter (no estaban XCharter ni newtxtext). Se compilaron
% y renderizaron a PNG cuatro candidatas —Pagella, Schola, Charter, Bonum— con
% el título de lámina y la banda de portada a tamaño real de proyección
% (comparar.tex/.pdf, carpeta de trabajo temporal). Resultado:
%   · TeX Gyre Schola y TeX Gyre Bonum: demasiado pesadas y de baja apertura,
%     compiten con la banda de color en vez de flotar sobre ella.
%   · Charter: elegante pero de contraste bajo y aire muy contemporáneo; se
%     aleja del carácter clásico-académico del modelo.
%   · TeX Gyre Pagella (clon de Palatino): contraste moderado, ojo grande,
%     aperturas abiertas —el «a» y el «g» de dos pisos que se ven en el
%     modelo—, y legible a tamaño de título incluso hasta el fondo de un
%     salón. Es además, junto con TeX Gyre Heros, parte de la misma familia
%     TeX Gyre: ambas fueron diseñadas para convivir, con métricas y peso
%     visual pensados para mezclarse. GANADORA.
%
% Se invoca la familia NFSS "qpl" directamente con \fontfamily, SIN cargar
% \usepackage{tgpagella}: ese paquete redefine \rmdefault a qpl, y aunque acá
% no se usa \rmfamily/\textrm/\mathrm en ninguna lámina (verificado con
% grep), no hay razón para tocar un valor global que no nos pertenece.
% \fontfamily{qpl} basta: NFSS localiza sola el .fd de TeX Gyre Pagella la
% primera vez que se usa. \familydefault SIGUE siendo \sfdefault: el cuerpo
% del deck no cambia.
%
% BUG CORREGIDO — \serifalab llevaba SOLO \fontfamily{qpl}, sin \selectfont.
% \fontfamily por sí solo no cambia la fuente activa: solo prepara el
% registro interno de NFSS: el cambio real ocurre hasta el siguiente
% \selectfont. En \setbeamerfont{title/subtitle/frametitle/framesubtitle}
% esto no se notaba porque la clave "family=" de beamer llama a \selectfont
% ella sola, al final de \usebeamerfont (ver beamerbasefont.sty). Pero en el
% pie de página \serifalab se usaba suelto, después de \usebeamerfont{footline}
% —que ya había hecho SU propio \selectfont en sans— y como nunca se volvía a
% llamar \selectfont, el cambio de familia quedaba preparado pero nunca se
% aplicaba: el pie seguía viéndose en sans siempre, sin importar cuántas veces
% se recompilara. \serifalab ahora incluye su propio \selectfont, así que
% funciona igual se use suelto o dentro de family=.
\newcommand{\serifalab}{\fontfamily{qpl}\selectfont}

% ── MATEMÁTICAS ───────────────────────────────────────────────────────────────
% Cargadas por si una lámina necesita un subíndice o una fracción.
% Recordatorio del §7: cero fórmulas con sumatorias en las láminas.
\usepackage{amsmath}
\usepackage{amssymb}

% ── GRÁFICAS Y TABLAS ─────────────────────────────────────────────────────────
\usepackage{graphicx}
\usepackage{booktabs}
\usepackage{tabularx}
\usepackage{array}

% ── COLOR Y CAJAS ─────────────────────────────────────────────────────────────
\usepackage{xcolor}
\input{../../estilo/paleta-itam.tex}   % ← única fuente de los hexes
\usepackage{tcolorbox}
\tcbuselibrary{skins, breakable}

\usepackage{tikz}
\usetikzlibrary{calc}

% ── CÓDIGO EN LAS LÁMINAS ─────────────────────────────────────────────────────
% El código va en las láminas porque muchas personas no alcanzan a copiar del
% proyector a su pantalla. Se usa listings y no minted: minted exige Python,
% Pygments y compilar con -shell-escape, y eso es fricción que nadie debería
% tener que reproducir para compilar el material de un curso.
\usepackage{listings}

\lstdefinestyle{Rlab}{
  language         = R,
  % \small y no \scriptsize: esto se proyecta en un salón y hay gente hasta el
  % fondo. Si un bloque no cabe a \small, el bloque es demasiado largo para una
  % lámina — se parte en dos, no se encoge la letra.
  basicstyle       = \ttfamily\small\color{grisTexto},
  keywordstyle     = \color{verdeITAM}\bfseries,
  % Los comentarios NO van en gris tenue: en este curso el comentario es el
  % contenido, no una nota al margen.
  commentstyle     = \color{grisTexto}\itshape,
  stringstyle      = \color{guinda},
  identifierstyle  = \color{grisTexto},
  backgroundcolor  = \color{verdeVelo},
  frame            = leftline,
  framesep         = 6pt,
  rulecolor        = \color{verdeITAM},
  framerule        = 1.2pt,
  xleftmargin      = 8pt,
  showstringspaces = false,
  breaklines       = true,
  breakatwhitespace= true,
  columns          = fullflexible,
  keepspaces       = true,
  upquote          = true,
  aboveskip        = 0.7em,
  belowskip        = 0.7em,
  literate =
    {á}{{\'a}}1 {é}{{\'e}}1 {í}{{\'i}}1 {ó}{{\'o}}1 {ú}{{\'u}}1
    {Á}{{\'A}}1 {É}{{\'E}}1 {Í}{{\'I}}1 {Ó}{{\'O}}1 {Ú}{{\'U}}1
    {ñ}{{\~n}}1 {Ñ}{{\~N}}1 {¿}{{?`}}1 {¡}{{!`}}1
    % El pipe nativo no es palabra reservada de R para listings y no puede
    % declararse como keyword porque no es alfanumérico. Se colorea a mano.
    % Es el operador que más aparece en todo el curso: tiene que verse.
    {|>}{{\textcolor{verdeITAM}{\bfseries|>}}}2
}
\lstset{style=Rlab}

% Verbos del tidyverse. No son palabras reservadas de R, pero en este curso son
% el vocabulario central: si se ven siempre del mismo color, el grupo empieza a
% reconocerlos antes de saber qué hacen.
\lstset{morekeywords={tibble, glimpse, mutate, filter, select, arrange,
  summarise, group_by, left_join, anti_join, count, ggplot, aes, labs,
  geom_col, geom_point, geom_sf, geom_histogram, st_read, st_crs, here, feols,
  pivot_longer, pivot_wider, read_csv, read_rds, clean_names, slice_max,
  case_when, theme_lab, ggsave, modelsummary, facet_wrap}}

% Código en línea dentro de un párrafo: \cod{filter()}
\newcommand{\cod}[1]{{\ttfamily\color{verdeProfundo}#1}}

% Énfasis. NUNCA \textbf en el cuerpo: el énfasis va por color.
\newcommand{\ac}[1]{\textcolor{verdeITAM}{#1}}   % acento primario
\newcommand{\se}[1]{\textcolor{guinda}{#1}}      % acento secundario, uso escaso

% ── LAYOUT Y UTILIDADES ───────────────────────────────────────────────────────
\usepackage{multicol}
\usepackage{appendixnumberbeamer}

% ── HIPERENLACES — siempre al final ───────────────────────────────────────────
\usepackage{url}
\usepackage{hyperref}

% ══════════════════════════════════════════════════════════════════════════════
% TEMA
% ══════════════════════════════════════════════════════════════════════════════
\usetheme{default}
\usecolortheme{default}
\usefonttheme{professionalfonts}   % NO {serif}: el deck es sans en todo
\setbeamertemplate{navigation symbols}{}

% ── Colores ───────────────────────────────────────────────────────────────────
\setbeamercolor{normal text}{fg=grisTexto, bg=white}
\setbeamercolor{structure}{fg=verdeITAM}

\setbeamercolor{frametitle}{fg=verdeITAM, bg=white}
\setbeamercolor{framesubtitle}{fg=grisTenue}

\setbeamercolor{item}{fg=verdeITAM}
\setbeamercolor{subitem}{fg=verdeITAM}
\setbeamercolor{subsubitem}{fg=grisTenue}

\setbeamercolor{block title}{fg=white, bg=verdeITAM}
\setbeamercolor{block body}{fg=grisTexto, bg=verdeVelo}
\setbeamercolor{block title alerted}{fg=white, bg=guinda}
\setbeamercolor{block body alerted}{fg=grisTexto, bg=grisVelo}
\setbeamercolor{block title example}{fg=verdeProfundo, bg=verdeTenue}
\setbeamercolor{block body example}{fg=grisTexto, bg=verdeVelo}

% Índice de las láminas de sección: la actual en blanco, las demás atenuadas.
% Se fijan SOLO los dos colores y se deja la plantilla de beamer intacta.
% Redefinir \setbeamertemplate{section in toc} a mano hace que la lámina de
% sección salga en blanco: probado, y no vale la pena averiguar por qué.
\setbeamercolor{section in toc}{fg=white}
\setbeamercolor{section in toc shaded}{fg=verdeSombra}

% ── Tamaños ───────────────────────────────────────────────────────────────────
% title, subtitle, frametitle y framesubtitle llevan family=\serifalab: son
% los cuatro elementos que "nombran" (portada y título de lámina). Todo lo
% demás —author, institute, date, block body, footline, section in toc— se
% deja SIN family, así que hereda \familydefault, que sigue siendo sans.
\setbeamerfont{title}{family=\serifalab, size=\LARGE}
\setbeamerfont{subtitle}{family=\serifalab, size=\large}
\setbeamerfont{author}{size=\normalsize}
\setbeamerfont{institute}{size=\small}
\setbeamerfont{date}{size=\small}
\setbeamerfont{frametitle}{family=\serifalab, size=\Large}
\setbeamerfont{framesubtitle}{family=\serifalab, size=\small}
\setbeamerfont{block title}{size=\normalsize, series=\bfseries}
\setbeamerfont{block body}{size=\small}
\setbeamerfont{footline}{size=\tiny}
\setbeamerfont{section in toc}{size=\large}

% ── Viñetas ───────────────────────────────────────────────────────────────────
% Punto medio pequeño, como en los decks de seminario. Nada de triángulos ni
% de cuadrados: la viñeta no debe competir con el texto.
\setbeamertemplate{itemize item}{\raisebox{0.12ex}{\scriptsize\textcolor{verdeITAM}{$\bullet$}}}
\setbeamertemplate{itemize subitem}{\raisebox{0.12ex}{\tiny\textcolor{verdeITAM}{$\bullet$}}}
\setbeamertemplate{enumerate item}{\textcolor{verdeITAM}{\insertenumlabel.}}
\setbeamertemplate{caption}{\insertcaption}

% ══════════════════════════════════════════════════════════════════════════════
% ENCABEZADO — banda superior con el logo del ITAM sobre verde
% ══════════════════════════════════════════════════════════════════════════════
\newcommand{\rutaLogo}{../../estilo/logo_itam.png}

\setbeamercolor{cajalogo}{fg=white, bg=verdeITAM}
\setbeamercolor{bandacurso}{fg=grisTenue, bg=white}

% Se arma con beamercolorbox y NO con tikz. Un `\fill` dentro de una
% tikzpicture con `overlay` en la plantilla de headline no se pinta —el nodo
% de la imagen sí aparece, pero el rectángulo de fondo no—, y como el
% lettering del ITAM es blanco sobre transparente, el logo quedaba invisible
% sobre el blanco de la lámina. Verificado con pdftoppm a 300 ppp.
% Historial: 2.4ex -> 3.0ex -> 3.5ex. Cada vez crecí la caja casi en la misma
% proporción que la imagen para evitar que se desbordara, y el resultado fue
% que la fracción de la caja verde que ocupa el logo casi no cambiaba (pasó de
% 43% a 46% a 49%): un cambio real pero demasiado sutil para notarse a
% simple vista. Esta vez el salto es deliberadamente grande y la caja crece
% MENOS que proporcionalmente, para que el logo se vea inconfundiblemente más
% grande y no solo "técnicamente" más grande: sube a 5.0ex (+43% sobre 3.5ex),
% con ht=5.2ex/dp=2.8ex (antes 4.6/2.6). El logo pasa a ocupar ~62% de la caja
% (antes ~49%). Sigue en el mismo lugar: esquina superior derecha, centrado.
% Recalculo r = (ht - dp)/2 - altura_imagen/2 = (5.2-2.8)/2 - 5.0/2 = -1.3ex.
% La caja izquierda lleva EXACTAMENTE los mismos ht y dp, o las dos mitades
% del encabezado quedan a distinta altura.
\setbeamertemplate{headline}{%
  \leavevmode%
  \hbox{%
    \begin{beamercolorbox}[wd=0.815\paperwidth, ht=5.2ex, dp=2.8ex, leftskip=1.4em]{bandacurso}%
      {\tiny LABORATORIO DE R PARA CIENCIA POLÍTICA}%
    \end{beamercolorbox}%
    % El logo se dimensiona en ex y no en fracción del ancho de página: así
    % cabe siempre dentro de la caja (ht 5.2ex + dp 2.8ex = 8.0ex de alto) y
    % no se desborda por arriba. El desplazamiento lo centra verticalmente.
    \begin{beamercolorbox}[wd=0.185\paperwidth, ht=5.2ex, dp=2.8ex, center]{cajalogo}%
      \raisebox{-1.3ex}{\includegraphics[height=5.0ex]{\rutaLogo}}%
    \end{beamercolorbox}%
  }%
  {\color{grisTexto!22}\hrule height 0.4pt}%
  \vskip0.7em%
}

% ══════════════════════════════════════════════════════════════════════════════
% PIE — sección a la izquierda, iniciales y folio a la derecha
% ══════════════════════════════════════════════════════════════════════════════
% Iniciales y folio ~1mm más a la izquierda (Cambio 3).
%
% DIAGNÓSTICO — antes de tocar el valor, se descubrió que "rightskip" en esta
% caja NUNCA había tenido efecto, ni con 1.4em ni con ningún otro valor: la
% clave "right" de beamercolorbox (beamerbasecolor.sty) hace
% \def\beamer@colbox@rs{0pt} y, como en la lista de opciones "right" iba
% DESPUÉS de "rightskip=1.4em", pisaba ese valor y lo dejaba en 0pt fijo.
% Medido en el PDF ya compilado: el «0» de «1/20» quedaba a 0.25mm del borde
% derecho de la lámina, no al 1.4em (≈5mm) que el código sugería. Es un bug
% latente de origen, no algo que haya introducido este cambio.
%
% ARREGLO — se invierte el orden ("right" primero, "rightskip" después) para
% que nuestro valor sea el que sobrevive. Y como el punto de partida real era
% 0pt (no 1.4em), para lograr el desplazamiento sutil de ~1mm que pidió EMG
% —"milimétricamente", no un recorte grande— el valor correcto NO es 2.2em
% (eso habría movido el bloque ~7.8mm, ocho veces más de lo pedido): es
% 0.45em, calibrado por render y medición en píxeles (antes: "0" de "1/20" a
% 0.25mm del borde; con 0.45em: a 1.02mm — ver reporte). La caja izquierda
% (sección) no se toca.
% Iniciales y folio («EMG | N/total») en serifa: \serifalab se aplica SOLO
% dentro de esta caja, después de \usebeamerfont{footline} (que solo trae el
% tamaño \tiny, sin family). El nombre de la sección de la caja izquierda
% NO lleva \serifalab y se queda en sans, como el resto del pie.
\setbeamertemplate{footline}{%
  \hbox{%
    \begin{beamercolorbox}[wd=0.62\paperwidth, ht=2.6ex, dp=1.4ex, leftskip=1.4em]{}%
      {\usebeamerfont{footline}\color{verdeITAM}\insertsectionhead}%
    \end{beamercolorbox}%
    \begin{beamercolorbox}[wd=0.38\paperwidth, ht=2.6ex, dp=1.4ex, right, rightskip=0.45em]{}%
      {\usebeamerfont{footline}\serifalab\color{grisTenue}%
       EMG\hspace{0.9em}\textbar\hspace{0.9em}%
       \insertframenumber{}\,/\,\inserttotalframenumber}%
    \end{beamercolorbox}%
  }%
  \vskip2pt%
}

% ══════════════════════════════════════════════════════════════════════════════
% TÍTULO DE FRAME — en verde, alineado a la izquierda, SIN regla debajo
% ══════════════════════════════════════════════════════════════════════════════
\setbeamertemplate{frametitle}{%
  \vspace{0.6em}
  \hspace{-0.4em}%
  \begin{minipage}{0.95\textwidth}
    \raggedright
    {\usebeamerfont{frametitle}\color{verdeITAM}\insertframetitle\par}%
    \ifx\insertframesubtitle\@empty\else%
      \vspace{0.15em}%
      {\usebeamerfont{framesubtitle}\color{grisTenue}\insertframesubtitle\par}%
    \fi
  \end{minipage}
  \vspace{0.5em}
}

% ══════════════════════════════════════════════════════════════════════════════
% PORTADA — banda verde con el título, datos debajo sobre blanco
% Título y subtítulo salen en serifa porque \usebeamerfont{title} y
% \usebeamerfont{subtitle} ya llevan family=\serifalab (ver §Tamaños): no hace
% falta tocar \inserttitle/\insertsubtitle aquí abajo. Autoría, adscripción y
% fecha usan \usebeamerfont{author/institute/date}, que NO llevan family y por
% lo tanto heredan \familydefault (sans) sin cambios.
% ══════════════════════════════════════════════════════════════════════════════
\setbeamertemplate{title page}{%
  \begin{tikzpicture}[remember picture, overlay]
    % Banda de color con el título dentro
    \fill[verdeITAM]
      ([yshift=-0.145\paperheight]current page.north west)
      rectangle ([yshift=-0.455\paperheight]current page.north east);
    \node[anchor=center, text width=0.86\paperwidth, align=center, inner sep=0pt]
      at ([yshift=-0.265\paperheight]current page.north)
      {\usebeamerfont{title}\color{white}\inserttitle};
    \node[anchor=center, text width=0.80\paperwidth, align=center, inner sep=0pt]
      at ([yshift=-0.385\paperheight]current page.north)
      {\usebeamerfont{subtitle}\color{verdeClaro}\insertsubtitle};

    % Autoría, adscripción y fecha
    \node[anchor=north, text width=0.86\paperwidth, align=center, inner sep=0pt]
      at ([yshift=-0.545\paperheight]current page.north)
      {\usebeamerfont{author}\color{grisTexto}\insertauthor};
    \node[anchor=north, text width=0.86\paperwidth, align=center, inner sep=0pt]
      at ([yshift=-0.625\paperheight]current page.north)
      {\usebeamerfont{institute}\color{grisTexto}\insertinstitute};
    \node[anchor=north, text width=0.86\paperwidth, align=center, inner sep=0pt]
      at ([yshift=-0.790\paperheight]current page.north)
      {\usebeamerfont{date}\color{grisTenue}\insertdate};
  \end{tikzpicture}%
}

% ══════════════════════════════════════════════════════════════════════════════
% LÁMINA DE SECCIÓN — fondo verde completo, índice con la sección actual en
% blanco y las demás atenuadas. Se dispara sola en cada \section{}.
% ══════════════════════════════════════════════════════════════════════════════
\AtBeginSection[]{%
  {%
  \setbeamercolor{background canvas}{bg=verdeITAM}
  % Las entradas del índice son hipervínculos, y con colorlinks=true hyperref
  % las pinta de linkcolor —que es el mismo verde del fondo—, así que la
  % sección actual desaparecía. Con hidelinks, hyperref deja de colorear y de
  % dibujar marcos, y entonces sí mandan los colores de beamer: la sección
  % actual en blanco y las demás atenuadas.
  % Diagnosticado a fuerza de bisecar el preámbulo; no es nada obvio.
  \hypersetup{hidelinks}
  \setbeamertemplate{headline}{}
  % Mismo arreglo y mismo valor que el footline normal (Cambio 3, ver
  % diagnóstico completo ahí arriba): "right" antes de "rightskip=0.45em", o la
  % clave "right" pisa nuestro valor y lo deja en 0pt. Si aquí se deja
  % distinto del footline normal, el pie de las láminas de sección queda
  % desalineado con el del resto.
  % Mismo \serifalab que el footline normal, y por la misma razón: solo en
  % «EMG | N/total», nunca en el nombre de la sección.
  \setbeamertemplate{footline}{%
    \hbox{%
      \begin{beamercolorbox}[wd=0.62\paperwidth, ht=2.6ex, dp=1.4ex, leftskip=1.4em]{}%
        {\usebeamerfont{footline}\color{white}\insertsectionhead}%
      \end{beamercolorbox}%
      \begin{beamercolorbox}[wd=0.38\paperwidth, ht=2.6ex, dp=1.4ex, right, rightskip=0.45em]{}%
        {\usebeamerfont{footline}\serifalab\color{verdeSombra}%
         EMG\hspace{0.9em}\textbar\hspace{0.9em}%
         \insertframenumber{}\,/\,\inserttotalframenumber}%
      \end{beamercolorbox}%
    }%
    \vskip2pt%
  }
  \begin{frame}[plain, noframenumbering]
    \vspace{1.4em}
    \begin{minipage}{0.86\textwidth}
      \hspace{1.2em}%
      \begin{minipage}{0.95\textwidth}
        \tableofcontents[currentsection, sectionstyle=show/shaded, hideallsubsections]
      \end{minipage}
    \end{minipage}
  \end{frame}
  }%
}

% ══════════════════════════════════════════════════════════════════════════════
% ENTORNOS PROPIOS DEL CURSO
% ══════════════════════════════════════════════════════════════════════════════

% Rótulo en versalitas falsas — tgheros sí tiene versalitas, pero el rótulo se
% ve mejor en mayúsculas con la letra pequeña.
\newcommand{\etiqueta}[1]{{\footnotesize\color{grisTenue}\MakeUppercase{#1}}}

% ── ¿Por qué nos importa esto? ────────────────────────────────────────────────
% Abre todas las sesiones: plantea la pregunta politológica antes de que
% aparezca una sola función.
\newcommand{\porquenosimporta}[2]{%
  \begin{frame}[plain, noframenumbering]
    \vfill
    \begin{center}
      \etiqueta{¿Por qué nos importa esto?}\\[1.3em]
      {\LARGE\color{verdeProfundo} #1}\\[1.5em]
      \begin{minipage}{0.80\textwidth}
        \centering\normalsize\color{grisTexto} #2
      \end{minipage}
    \end{center}
    \vfill
  \end{frame}%
}

% ── Lo que ya sabemos hacer ───────────────────────────────────────────────────
% Cierra todas las sesiones. Exactamente tres viñetas: es la recompensa.
\newcommand{\yasabemos}[3]{%
  \begin{frame}{Lo que ya sabemos hacer}
    \vspace{1.0em}
    \begin{itemize}\setlength{\itemsep}{1.4em}
      \item #1
      \item #2
      \item #3
    \end{itemize}
    \vfill
  \end{frame}%
}

% ── Adelanto de la sesión siguiente ───────────────────────────────────────────
\newcommand{\lasiguiente}[2]{%
  \begin{frame}[plain, noframenumbering]
    \vfill
    \begin{center}
      \etiqueta{La semana que viene}\\[1.3em]
      {\LARGE\color{verdeITAM} #1}\\[1.5em]
      \begin{minipage}{0.80\textwidth}
        \centering\normalsize\color{grisTexto} #2
      \end{minipage}
    \end{center}
    \vfill
  \end{frame}%
}

% ── Caja de error anticipado ──────────────────────────────────────────────────
% Los errores son parte del temario. Esta caja los presenta ANTES de que
% ocurran, con su traducción del inglés y su arreglo.
% "no shadow" y "sharp corners" son deliberados: nada de sombras ni degradados.
\newtcolorbox{errorbox}[1]{
  enhanced, breakable, no shadow, sharp corners,
  colback   = grisVelo,
  colframe  = guinda,
  boxrule   = 0pt, leftrule = 2.5pt,
  toprule = 0pt, bottomrule = 0pt, rightrule = 0pt,
  fontupper = \small\color{grisTexto},
  title     = {\color{guinda}\bfseries\small\ttfamily #1},
  colbacktitle = grisVelo,
  coltitle  = guinda,
  titlerule = 0pt
}

% ── Caja del salto politólogo ─────────────────────────────────────────────────
% El §7 pide que cada sesión contenga al menos un momento donde la intuición
% política se convierta en pregunta empírica, y que ese momento se nombre.
% Esta caja es el lugar donde se nombra.
\newtcolorbox{saltobox}{
  enhanced, breakable, no shadow, sharp corners,
  colback   = verdeVelo,
  colframe  = verdeITAM,
  boxrule   = 0pt, leftrule = 2.5pt,
  toprule = 0pt, bottomrule = 0pt, rightrule = 0pt,
  fontupper = \small\color{grisTexto},
  title     = {\color{verdeProfundo}\bfseries\small El salto},
  colbacktitle = verdeVelo,
  coltitle  = verdeProfundo,
  titlerule = 0pt
}

% ══════════════════════════════════════════════════════════════════════════════
% METADATOS FIJOS
% ══════════════════════════════════════════════════════════════════════════════
\author[E. Miranda]{Emiliano Miranda González}
\institute[ITAM]{%
  Instituto Tecnológico Autónomo de México\\
  Departamento Académico de Ciencia Política%
}

\hypersetup{
  colorlinks  = true,
  linkcolor   = verdeITAM,
  urlcolor    = verdeITAM,
  citecolor   = grisTenue,
  pdfauthor   = {Emiliano Miranda González},
  pdfsubject  = {Laboratorio de R para Ciencia Política — ITAM},
  pdfkeywords = {R, ciencia política, ITAM, métodos cuantitativos},
  unicode     = true
}

%
% Se compila con pdflatex, DOS VECES (los nodos de tikz con `remember picture`
% necesitan la segunda pasada para colocarse bien).
% ══════════════════════════════════════════════════════════════════════════════

\documentclass[
  aspectratio=43,
  t,
  10pt
]{beamer}

% ── TIPOGRAFÍA ────────────────────────────────────────────────────────────────
% El cuerpo del deck sigue siendo sans en su totalidad. TeX Gyre Heros es un
% clon de Helvetica en Type 1 con cobertura T1 completa —incluidas cursivas y
% versalitas reales—, presente en cualquier distribución de TeX Live. lmodern
% va ANTES para quedarse con la monoespaciada: sin él, los comentarios en
% cursiva de los bloques de código salen como fuente de mapa de bits y se ven
% borrosos al proyectar.
%
% Si en la máquina de quien compila están instaladas Fira Sans o Lato —que son
% las humanistas que suelen usar los decks de seminario de ciencia política—,
% basta sustituir la línea de tgheros por \usepackage[sfdefault]{FiraSans} o
% \usepackage[default]{lato}. El resto del preámbulo no cambia.
\usepackage[T1]{fontenc}
\usepackage[utf8]{inputenc}
\usepackage{lmodern}
\usepackage{tgheros}
\renewcommand{\familydefault}{\sfdefault}
\usepackage{microtype}

% ── SERIFA PARA TÍTULOS (mezcla del deck de seminario) ─────────────────────────
% EMG pidió la mezcla de los decks de seminario de ciencia política
% internacional (ref.: deck de Gary King, IQSS/Harvard): serifa para lo que
% NOMBRA —portada y título de cada lámina— y sans para lo que INFORMA —cuerpo,
% viñetas, índice de secciones y pie—. Ver \serifalab más abajo y su uso en
% \setbeamerfont{title/subtitle/frametitle/framesubtitle}.
%
% Elección de la familia, verificada y no adivinada. De las candidatas del
% encargo, kpsewhich confirmó instaladas: tgpagella, tgtermes, tgschola,
% tgbonum, mathpazo, charter (no estaban XCharter ni newtxtext). Se compilaron
% y renderizaron a PNG cuatro candidatas —Pagella, Schola, Charter, Bonum— con
% el título de lámina y la banda de portada a tamaño real de proyección
% (comparar.tex/.pdf, carpeta de trabajo temporal). Resultado:
%   · TeX Gyre Schola y TeX Gyre Bonum: demasiado pesadas y de baja apertura,
%     compiten con la banda de color en vez de flotar sobre ella.
%   · Charter: elegante pero de contraste bajo y aire muy contemporáneo; se
%     aleja del carácter clásico-académico del modelo.
%   · TeX Gyre Pagella (clon de Palatino): contraste moderado, ojo grande,
%     aperturas abiertas —el «a» y el «g» de dos pisos que se ven en el
%     modelo—, y legible a tamaño de título incluso hasta el fondo de un
%     salón. Es además, junto con TeX Gyre Heros, parte de la misma familia
%     TeX Gyre: ambas fueron diseñadas para convivir, con métricas y peso
%     visual pensados para mezclarse. GANADORA.
%
% Se invoca la familia NFSS "qpl" directamente con \fontfamily, SIN cargar
% \usepackage{tgpagella}: ese paquete redefine \rmdefault a qpl, y aunque acá
% no se usa \rmfamily/\textrm/\mathrm en ninguna lámina (verificado con
% grep), no hay razón para tocar un valor global que no nos pertenece.
% \fontfamily{qpl} basta: NFSS localiza sola el .fd de TeX Gyre Pagella la
% primera vez que se usa. \familydefault SIGUE siendo \sfdefault: el cuerpo
% del deck no cambia.
%
% BUG CORREGIDO — \serifalab llevaba SOLO \fontfamily{qpl}, sin \selectfont.
% \fontfamily por sí solo no cambia la fuente activa: solo prepara el
% registro interno de NFSS: el cambio real ocurre hasta el siguiente
% \selectfont. En \setbeamerfont{title/subtitle/frametitle/framesubtitle}
% esto no se notaba porque la clave "family=" de beamer llama a \selectfont
% ella sola, al final de \usebeamerfont (ver beamerbasefont.sty). Pero en el
% pie de página \serifalab se usaba suelto, después de \usebeamerfont{footline}
% —que ya había hecho SU propio \selectfont en sans— y como nunca se volvía a
% llamar \selectfont, el cambio de familia quedaba preparado pero nunca se
% aplicaba: el pie seguía viéndose en sans siempre, sin importar cuántas veces
% se recompilara. \serifalab ahora incluye su propio \selectfont, así que
% funciona igual se use suelto o dentro de family=.
\newcommand{\serifalab}{\fontfamily{qpl}\selectfont}

% ── MATEMÁTICAS ───────────────────────────────────────────────────────────────
% Cargadas por si una lámina necesita un subíndice o una fracción.
% Recordatorio del §7: cero fórmulas con sumatorias en las láminas.
\usepackage{amsmath}
\usepackage{amssymb}

% ── GRÁFICAS Y TABLAS ─────────────────────────────────────────────────────────
\usepackage{graphicx}
\usepackage{booktabs}
\usepackage{tabularx}
\usepackage{array}

% ── COLOR Y CAJAS ─────────────────────────────────────────────────────────────
\usepackage{xcolor}
% ══════════════════════════════════════════════════════════════════════════════
% PALETA DEL LABORATORIO DE R PARA CIENCIA POLÍTICA — ITAM
% Verde institucional #006847 sobre blanco
%
% ESTE ES EL ÚNICO LUGAR NORMATIVO DONDE VIVEN LOS HEXES.
% Si cambia un color, se cambia aquí y se propaga a mano a los otros dos
% archivos que lo heredan:
%   · estilo/tema_lab.R     → theme_lab() y las escalas de ggplot2
%   · estilo/estilo.scss    → el sitio Quarto
% Los tres tienen que decir lo mismo. Si no, las láminas, los gráficos y el
% sitio dejan de combinar y se nota.
% ══════════════════════════════════════════════════════════════════════════════

\definecolor{verdeITAM}{HTML}{006847}     % acento principal — verde institucional
\definecolor{verdeProfundo}{HTML}{00402C} % títulos y texto enfático
\definecolor{verdeClaro}{HTML}{A8D5BE}    % fondos suaves, resaltados
\definecolor{grisTexto}{HTML}{3D3D3D}     % cuerpo de texto
\definecolor{guinda}{HTML}{7B1113}        % acento secundario, uso escaso

% ── Derivados ────────────────────────────────────────────────────────────────
% No son colores nuevos: son diluciones de los anteriores. Se definen aquí para
% que nadie escriba "verdeClaro!15" a mano en una lámina y el curso termine con
% quince tonos distintos repartidos por ahí.

\colorlet{verdeVelo}{verdeClaro!22!white}   % fondo de bloque y de código
\colorlet{verdeTenue}{verdeClaro!55!white}  % resaltado dentro de un bloque
\colorlet{verdeSombra}{white!30!verdeITAM}  % texto atenuado SOBRE fondo verde
\colorlet{grisVelo}{grisTexto!7!white}      % fondos muy tenues, reglas
\colorlet{grisTenue}{grisTexto!45!white}    % texto terciario, pies
   % ← única fuente de los hexes
\usepackage{tcolorbox}
\tcbuselibrary{skins, breakable}

\usepackage{tikz}
\usetikzlibrary{calc}

% ── CÓDIGO EN LAS LÁMINAS ─────────────────────────────────────────────────────
% El código va en las láminas porque muchas personas no alcanzan a copiar del
% proyector a su pantalla. Se usa listings y no minted: minted exige Python,
% Pygments y compilar con -shell-escape, y eso es fricción que nadie debería
% tener que reproducir para compilar el material de un curso.
\usepackage{listings}

\lstdefinestyle{Rlab}{
  language         = R,
  % \small y no \scriptsize: esto se proyecta en un salón y hay gente hasta el
  % fondo. Si un bloque no cabe a \small, el bloque es demasiado largo para una
  % lámina — se parte en dos, no se encoge la letra.
  basicstyle       = \ttfamily\small\color{grisTexto},
  keywordstyle     = \color{verdeITAM}\bfseries,
  % Los comentarios NO van en gris tenue: en este curso el comentario es el
  % contenido, no una nota al margen.
  commentstyle     = \color{grisTexto}\itshape,
  stringstyle      = \color{guinda},
  identifierstyle  = \color{grisTexto},
  backgroundcolor  = \color{verdeVelo},
  frame            = leftline,
  framesep         = 6pt,
  rulecolor        = \color{verdeITAM},
  framerule        = 1.2pt,
  xleftmargin      = 8pt,
  showstringspaces = false,
  breaklines       = true,
  breakatwhitespace= true,
  columns          = fullflexible,
  keepspaces       = true,
  upquote          = true,
  aboveskip        = 0.7em,
  belowskip        = 0.7em,
  literate =
    {á}{{\'a}}1 {é}{{\'e}}1 {í}{{\'i}}1 {ó}{{\'o}}1 {ú}{{\'u}}1
    {Á}{{\'A}}1 {É}{{\'E}}1 {Í}{{\'I}}1 {Ó}{{\'O}}1 {Ú}{{\'U}}1
    {ñ}{{\~n}}1 {Ñ}{{\~N}}1 {¿}{{?`}}1 {¡}{{!`}}1
    % El pipe nativo no es palabra reservada de R para listings y no puede
    % declararse como keyword porque no es alfanumérico. Se colorea a mano.
    % Es el operador que más aparece en todo el curso: tiene que verse.
    {|>}{{\textcolor{verdeITAM}{\bfseries|>}}}2
}
\lstset{style=Rlab}

% Verbos del tidyverse. No son palabras reservadas de R, pero en este curso son
% el vocabulario central: si se ven siempre del mismo color, el grupo empieza a
% reconocerlos antes de saber qué hacen.
\lstset{morekeywords={tibble, glimpse, mutate, filter, select, arrange,
  summarise, group_by, left_join, anti_join, count, ggplot, aes, labs,
  geom_col, geom_point, geom_sf, geom_histogram, st_read, st_crs, here, feols,
  pivot_longer, pivot_wider, read_csv, read_rds, clean_names, slice_max,
  case_when, theme_lab, ggsave, modelsummary, facet_wrap}}

% Código en línea dentro de un párrafo: \cod{filter()}
\newcommand{\cod}[1]{{\ttfamily\color{verdeProfundo}#1}}

% Énfasis. NUNCA \textbf en el cuerpo: el énfasis va por color.
\newcommand{\ac}[1]{\textcolor{verdeITAM}{#1}}   % acento primario
\newcommand{\se}[1]{\textcolor{guinda}{#1}}      % acento secundario, uso escaso

% ── LAYOUT Y UTILIDADES ───────────────────────────────────────────────────────
\usepackage{multicol}
\usepackage{appendixnumberbeamer}

% ── HIPERENLACES — siempre al final ───────────────────────────────────────────
\usepackage{url}
\usepackage{hyperref}

% ══════════════════════════════════════════════════════════════════════════════
% TEMA
% ══════════════════════════════════════════════════════════════════════════════
\usetheme{default}
\usecolortheme{default}
\usefonttheme{professionalfonts}   % NO {serif}: el deck es sans en todo
\setbeamertemplate{navigation symbols}{}

% ── Colores ───────────────────────────────────────────────────────────────────
\setbeamercolor{normal text}{fg=grisTexto, bg=white}
\setbeamercolor{structure}{fg=verdeITAM}

\setbeamercolor{frametitle}{fg=verdeITAM, bg=white}
\setbeamercolor{framesubtitle}{fg=grisTenue}

\setbeamercolor{item}{fg=verdeITAM}
\setbeamercolor{subitem}{fg=verdeITAM}
\setbeamercolor{subsubitem}{fg=grisTenue}

\setbeamercolor{block title}{fg=white, bg=verdeITAM}
\setbeamercolor{block body}{fg=grisTexto, bg=verdeVelo}
\setbeamercolor{block title alerted}{fg=white, bg=guinda}
\setbeamercolor{block body alerted}{fg=grisTexto, bg=grisVelo}
\setbeamercolor{block title example}{fg=verdeProfundo, bg=verdeTenue}
\setbeamercolor{block body example}{fg=grisTexto, bg=verdeVelo}

% Índice de las láminas de sección: la actual en blanco, las demás atenuadas.
% Se fijan SOLO los dos colores y se deja la plantilla de beamer intacta.
% Redefinir \setbeamertemplate{section in toc} a mano hace que la lámina de
% sección salga en blanco: probado, y no vale la pena averiguar por qué.
\setbeamercolor{section in toc}{fg=white}
\setbeamercolor{section in toc shaded}{fg=verdeSombra}

% ── Tamaños ───────────────────────────────────────────────────────────────────
% title, subtitle, frametitle y framesubtitle llevan family=\serifalab: son
% los cuatro elementos que "nombran" (portada y título de lámina). Todo lo
% demás —author, institute, date, block body, footline, section in toc— se
% deja SIN family, así que hereda \familydefault, que sigue siendo sans.
\setbeamerfont{title}{family=\serifalab, size=\LARGE}
\setbeamerfont{subtitle}{family=\serifalab, size=\large}
\setbeamerfont{author}{size=\normalsize}
\setbeamerfont{institute}{size=\small}
\setbeamerfont{date}{size=\small}
\setbeamerfont{frametitle}{family=\serifalab, size=\Large}
\setbeamerfont{framesubtitle}{family=\serifalab, size=\small}
\setbeamerfont{block title}{size=\normalsize, series=\bfseries}
\setbeamerfont{block body}{size=\small}
\setbeamerfont{footline}{size=\tiny}
\setbeamerfont{section in toc}{size=\large}

% ── Viñetas ───────────────────────────────────────────────────────────────────
% Punto medio pequeño, como en los decks de seminario. Nada de triángulos ni
% de cuadrados: la viñeta no debe competir con el texto.
\setbeamertemplate{itemize item}{\raisebox{0.12ex}{\scriptsize\textcolor{verdeITAM}{$\bullet$}}}
\setbeamertemplate{itemize subitem}{\raisebox{0.12ex}{\tiny\textcolor{verdeITAM}{$\bullet$}}}
\setbeamertemplate{enumerate item}{\textcolor{verdeITAM}{\insertenumlabel.}}
\setbeamertemplate{caption}{\insertcaption}

% ══════════════════════════════════════════════════════════════════════════════
% ENCABEZADO — banda superior con el logo del ITAM sobre verde
% ══════════════════════════════════════════════════════════════════════════════
\newcommand{\rutaLogo}{../../estilo/logo_itam.png}

\setbeamercolor{cajalogo}{fg=white, bg=verdeITAM}
\setbeamercolor{bandacurso}{fg=grisTenue, bg=white}

% Se arma con beamercolorbox y NO con tikz. Un `\fill` dentro de una
% tikzpicture con `overlay` en la plantilla de headline no se pinta —el nodo
% de la imagen sí aparece, pero el rectángulo de fondo no—, y como el
% lettering del ITAM es blanco sobre transparente, el logo quedaba invisible
% sobre el blanco de la lámina. Verificado con pdftoppm a 300 ppp.
% Historial: 2.4ex -> 3.0ex -> 3.5ex. Cada vez crecí la caja casi en la misma
% proporción que la imagen para evitar que se desbordara, y el resultado fue
% que la fracción de la caja verde que ocupa el logo casi no cambiaba (pasó de
% 43% a 46% a 49%): un cambio real pero demasiado sutil para notarse a
% simple vista. Esta vez el salto es deliberadamente grande y la caja crece
% MENOS que proporcionalmente, para que el logo se vea inconfundiblemente más
% grande y no solo "técnicamente" más grande: sube a 5.0ex (+43% sobre 3.5ex),
% con ht=5.2ex/dp=2.8ex (antes 4.6/2.6). El logo pasa a ocupar ~62% de la caja
% (antes ~49%). Sigue en el mismo lugar: esquina superior derecha, centrado.
% Recalculo r = (ht - dp)/2 - altura_imagen/2 = (5.2-2.8)/2 - 5.0/2 = -1.3ex.
% La caja izquierda lleva EXACTAMENTE los mismos ht y dp, o las dos mitades
% del encabezado quedan a distinta altura.
\setbeamertemplate{headline}{%
  \leavevmode%
  \hbox{%
    \begin{beamercolorbox}[wd=0.815\paperwidth, ht=5.2ex, dp=2.8ex, leftskip=1.4em]{bandacurso}%
      {\tiny LABORATORIO DE R PARA CIENCIA POLÍTICA}%
    \end{beamercolorbox}%
    % El logo se dimensiona en ex y no en fracción del ancho de página: así
    % cabe siempre dentro de la caja (ht 5.2ex + dp 2.8ex = 8.0ex de alto) y
    % no se desborda por arriba. El desplazamiento lo centra verticalmente.
    \begin{beamercolorbox}[wd=0.185\paperwidth, ht=5.2ex, dp=2.8ex, center]{cajalogo}%
      \raisebox{-1.3ex}{\includegraphics[height=5.0ex]{\rutaLogo}}%
    \end{beamercolorbox}%
  }%
  {\color{grisTexto!22}\hrule height 0.4pt}%
  \vskip0.7em%
}

% ══════════════════════════════════════════════════════════════════════════════
% PIE — sección a la izquierda, iniciales y folio a la derecha
% ══════════════════════════════════════════════════════════════════════════════
% Iniciales y folio ~1mm más a la izquierda (Cambio 3).
%
% DIAGNÓSTICO — antes de tocar el valor, se descubrió que "rightskip" en esta
% caja NUNCA había tenido efecto, ni con 1.4em ni con ningún otro valor: la
% clave "right" de beamercolorbox (beamerbasecolor.sty) hace
% \def\beamer@colbox@rs{0pt} y, como en la lista de opciones "right" iba
% DESPUÉS de "rightskip=1.4em", pisaba ese valor y lo dejaba en 0pt fijo.
% Medido en el PDF ya compilado: el «0» de «1/20» quedaba a 0.25mm del borde
% derecho de la lámina, no al 1.4em (≈5mm) que el código sugería. Es un bug
% latente de origen, no algo que haya introducido este cambio.
%
% ARREGLO — se invierte el orden ("right" primero, "rightskip" después) para
% que nuestro valor sea el que sobrevive. Y como el punto de partida real era
% 0pt (no 1.4em), para lograr el desplazamiento sutil de ~1mm que pidió EMG
% —"milimétricamente", no un recorte grande— el valor correcto NO es 2.2em
% (eso habría movido el bloque ~7.8mm, ocho veces más de lo pedido): es
% 0.45em, calibrado por render y medición en píxeles (antes: "0" de "1/20" a
% 0.25mm del borde; con 0.45em: a 1.02mm — ver reporte). La caja izquierda
% (sección) no se toca.
% Iniciales y folio («EMG | N/total») en serifa: \serifalab se aplica SOLO
% dentro de esta caja, después de \usebeamerfont{footline} (que solo trae el
% tamaño \tiny, sin family). El nombre de la sección de la caja izquierda
% NO lleva \serifalab y se queda en sans, como el resto del pie.
\setbeamertemplate{footline}{%
  \hbox{%
    \begin{beamercolorbox}[wd=0.62\paperwidth, ht=2.6ex, dp=1.4ex, leftskip=1.4em]{}%
      {\usebeamerfont{footline}\color{verdeITAM}\insertsectionhead}%
    \end{beamercolorbox}%
    \begin{beamercolorbox}[wd=0.38\paperwidth, ht=2.6ex, dp=1.4ex, right, rightskip=0.45em]{}%
      {\usebeamerfont{footline}\serifalab\color{grisTenue}%
       EMG\hspace{0.9em}\textbar\hspace{0.9em}%
       \insertframenumber{}\,/\,\inserttotalframenumber}%
    \end{beamercolorbox}%
  }%
  \vskip2pt%
}

% ══════════════════════════════════════════════════════════════════════════════
% TÍTULO DE FRAME — en verde, alineado a la izquierda, SIN regla debajo
% ══════════════════════════════════════════════════════════════════════════════
\setbeamertemplate{frametitle}{%
  \vspace{0.6em}
  \hspace{-0.4em}%
  \begin{minipage}{0.95\textwidth}
    \raggedright
    {\usebeamerfont{frametitle}\color{verdeITAM}\insertframetitle\par}%
    \ifx\insertframesubtitle\@empty\else%
      \vspace{0.15em}%
      {\usebeamerfont{framesubtitle}\color{grisTenue}\insertframesubtitle\par}%
    \fi
  \end{minipage}
  \vspace{0.5em}
}

% ══════════════════════════════════════════════════════════════════════════════
% PORTADA — banda verde con el título, datos debajo sobre blanco
% Título y subtítulo salen en serifa porque \usebeamerfont{title} y
% \usebeamerfont{subtitle} ya llevan family=\serifalab (ver §Tamaños): no hace
% falta tocar \inserttitle/\insertsubtitle aquí abajo. Autoría, adscripción y
% fecha usan \usebeamerfont{author/institute/date}, que NO llevan family y por
% lo tanto heredan \familydefault (sans) sin cambios.
% ══════════════════════════════════════════════════════════════════════════════
\setbeamertemplate{title page}{%
  \begin{tikzpicture}[remember picture, overlay]
    % Banda de color con el título dentro
    \fill[verdeITAM]
      ([yshift=-0.145\paperheight]current page.north west)
      rectangle ([yshift=-0.455\paperheight]current page.north east);
    \node[anchor=center, text width=0.86\paperwidth, align=center, inner sep=0pt]
      at ([yshift=-0.265\paperheight]current page.north)
      {\usebeamerfont{title}\color{white}\inserttitle};
    \node[anchor=center, text width=0.80\paperwidth, align=center, inner sep=0pt]
      at ([yshift=-0.385\paperheight]current page.north)
      {\usebeamerfont{subtitle}\color{verdeClaro}\insertsubtitle};

    % Autoría, adscripción y fecha
    \node[anchor=north, text width=0.86\paperwidth, align=center, inner sep=0pt]
      at ([yshift=-0.545\paperheight]current page.north)
      {\usebeamerfont{author}\color{grisTexto}\insertauthor};
    \node[anchor=north, text width=0.86\paperwidth, align=center, inner sep=0pt]
      at ([yshift=-0.625\paperheight]current page.north)
      {\usebeamerfont{institute}\color{grisTexto}\insertinstitute};
    \node[anchor=north, text width=0.86\paperwidth, align=center, inner sep=0pt]
      at ([yshift=-0.790\paperheight]current page.north)
      {\usebeamerfont{date}\color{grisTenue}\insertdate};
  \end{tikzpicture}%
}

% ══════════════════════════════════════════════════════════════════════════════
% LÁMINA DE SECCIÓN — fondo verde completo, índice con la sección actual en
% blanco y las demás atenuadas. Se dispara sola en cada \section{}.
% ══════════════════════════════════════════════════════════════════════════════
\AtBeginSection[]{%
  {%
  \setbeamercolor{background canvas}{bg=verdeITAM}
  % Las entradas del índice son hipervínculos, y con colorlinks=true hyperref
  % las pinta de linkcolor —que es el mismo verde del fondo—, así que la
  % sección actual desaparecía. Con hidelinks, hyperref deja de colorear y de
  % dibujar marcos, y entonces sí mandan los colores de beamer: la sección
  % actual en blanco y las demás atenuadas.
  % Diagnosticado a fuerza de bisecar el preámbulo; no es nada obvio.
  \hypersetup{hidelinks}
  \setbeamertemplate{headline}{}
  % Mismo arreglo y mismo valor que el footline normal (Cambio 3, ver
  % diagnóstico completo ahí arriba): "right" antes de "rightskip=0.45em", o la
  % clave "right" pisa nuestro valor y lo deja en 0pt. Si aquí se deja
  % distinto del footline normal, el pie de las láminas de sección queda
  % desalineado con el del resto.
  % Mismo \serifalab que el footline normal, y por la misma razón: solo en
  % «EMG | N/total», nunca en el nombre de la sección.
  \setbeamertemplate{footline}{%
    \hbox{%
      \begin{beamercolorbox}[wd=0.62\paperwidth, ht=2.6ex, dp=1.4ex, leftskip=1.4em]{}%
        {\usebeamerfont{footline}\color{white}\insertsectionhead}%
      \end{beamercolorbox}%
      \begin{beamercolorbox}[wd=0.38\paperwidth, ht=2.6ex, dp=1.4ex, right, rightskip=0.45em]{}%
        {\usebeamerfont{footline}\serifalab\color{verdeSombra}%
         EMG\hspace{0.9em}\textbar\hspace{0.9em}%
         \insertframenumber{}\,/\,\inserttotalframenumber}%
      \end{beamercolorbox}%
    }%
    \vskip2pt%
  }
  \begin{frame}[plain, noframenumbering]
    \vspace{1.4em}
    \begin{minipage}{0.86\textwidth}
      \hspace{1.2em}%
      \begin{minipage}{0.95\textwidth}
        \tableofcontents[currentsection, sectionstyle=show/shaded, hideallsubsections]
      \end{minipage}
    \end{minipage}
  \end{frame}
  }%
}

% ══════════════════════════════════════════════════════════════════════════════
% ENTORNOS PROPIOS DEL CURSO
% ══════════════════════════════════════════════════════════════════════════════

% Rótulo en versalitas falsas — tgheros sí tiene versalitas, pero el rótulo se
% ve mejor en mayúsculas con la letra pequeña.
\newcommand{\etiqueta}[1]{{\footnotesize\color{grisTenue}\MakeUppercase{#1}}}

% ── ¿Por qué nos importa esto? ────────────────────────────────────────────────
% Abre todas las sesiones: plantea la pregunta politológica antes de que
% aparezca una sola función.
\newcommand{\porquenosimporta}[2]{%
  \begin{frame}[plain, noframenumbering]
    \vfill
    \begin{center}
      \etiqueta{¿Por qué nos importa esto?}\\[1.3em]
      {\LARGE\color{verdeProfundo} #1}\\[1.5em]
      \begin{minipage}{0.80\textwidth}
        \centering\normalsize\color{grisTexto} #2
      \end{minipage}
    \end{center}
    \vfill
  \end{frame}%
}

% ── Lo que ya sabemos hacer ───────────────────────────────────────────────────
% Cierra todas las sesiones. Exactamente tres viñetas: es la recompensa.
\newcommand{\yasabemos}[3]{%
  \begin{frame}{Lo que ya sabemos hacer}
    \vspace{1.0em}
    \begin{itemize}\setlength{\itemsep}{1.4em}
      \item #1
      \item #2
      \item #3
    \end{itemize}
    \vfill
  \end{frame}%
}

% ── Adelanto de la sesión siguiente ───────────────────────────────────────────
\newcommand{\lasiguiente}[2]{%
  \begin{frame}[plain, noframenumbering]
    \vfill
    \begin{center}
      \etiqueta{La semana que viene}\\[1.3em]
      {\LARGE\color{verdeITAM} #1}\\[1.5em]
      \begin{minipage}{0.80\textwidth}
        \centering\normalsize\color{grisTexto} #2
      \end{minipage}
    \end{center}
    \vfill
  \end{frame}%
}

% ── Caja de error anticipado ──────────────────────────────────────────────────
% Los errores son parte del temario. Esta caja los presenta ANTES de que
% ocurran, con su traducción del inglés y su arreglo.
% "no shadow" y "sharp corners" son deliberados: nada de sombras ni degradados.
\newtcolorbox{errorbox}[1]{
  enhanced, breakable, no shadow, sharp corners,
  colback   = grisVelo,
  colframe  = guinda,
  boxrule   = 0pt, leftrule = 2.5pt,
  toprule = 0pt, bottomrule = 0pt, rightrule = 0pt,
  fontupper = \small\color{grisTexto},
  title     = {\color{guinda}\bfseries\small\ttfamily #1},
  colbacktitle = grisVelo,
  coltitle  = guinda,
  titlerule = 0pt
}

% ── Caja del salto politólogo ─────────────────────────────────────────────────
% El §7 pide que cada sesión contenga al menos un momento donde la intuición
% política se convierta en pregunta empírica, y que ese momento se nombre.
% Esta caja es el lugar donde se nombra.
\newtcolorbox{saltobox}{
  enhanced, breakable, no shadow, sharp corners,
  colback   = verdeVelo,
  colframe  = verdeITAM,
  boxrule   = 0pt, leftrule = 2.5pt,
  toprule = 0pt, bottomrule = 0pt, rightrule = 0pt,
  fontupper = \small\color{grisTexto},
  title     = {\color{verdeProfundo}\bfseries\small El salto},
  colbacktitle = verdeVelo,
  coltitle  = verdeProfundo,
  titlerule = 0pt
}

% ══════════════════════════════════════════════════════════════════════════════
% METADATOS FIJOS
% ══════════════════════════════════════════════════════════════════════════════
\author[E. Miranda]{Emiliano Miranda González}
\institute[ITAM]{%
  Instituto Tecnológico Autónomo de México\\
  Departamento Académico de Ciencia Política%
}

\hypersetup{
  colorlinks  = true,
  linkcolor   = verdeITAM,
  urlcolor    = verdeITAM,
  citecolor   = grisTenue,
  pdfauthor   = {Emiliano Miranda González},
  pdfsubject  = {Laboratorio de R para Ciencia Política — ITAM},
  pdfkeywords = {R, ciencia política, ITAM, métodos cuantitativos},
  unicode     = true
}

%
% Se compila con pdflatex, DOS VECES (los nodos de tikz con `remember picture`
% necesitan la segunda pasada para colocarse bien).
% ══════════════════════════════════════════════════════════════════════════════

\documentclass[
  aspectratio=43,
  t,
  10pt
]{beamer}

% ── TIPOGRAFÍA ────────────────────────────────────────────────────────────────
% El cuerpo del deck sigue siendo sans en su totalidad. TeX Gyre Heros es un
% clon de Helvetica en Type 1 con cobertura T1 completa —incluidas cursivas y
% versalitas reales—, presente en cualquier distribución de TeX Live. lmodern
% va ANTES para quedarse con la monoespaciada: sin él, los comentarios en
% cursiva de los bloques de código salen como fuente de mapa de bits y se ven
% borrosos al proyectar.
%
% Si en la máquina de quien compila están instaladas Fira Sans o Lato —que son
% las humanistas que suelen usar los decks de seminario de ciencia política—,
% basta sustituir la línea de tgheros por \usepackage[sfdefault]{FiraSans} o
% \usepackage[default]{lato}. El resto del preámbulo no cambia.
\usepackage[T1]{fontenc}
\usepackage[utf8]{inputenc}
\usepackage{lmodern}
\usepackage{tgheros}
\renewcommand{\familydefault}{\sfdefault}
\usepackage{microtype}

% ── SERIFA PARA TÍTULOS (mezcla del deck de seminario) ─────────────────────────
% EMG pidió la mezcla de los decks de seminario de ciencia política
% internacional (ref.: deck de Gary King, IQSS/Harvard): serifa para lo que
% NOMBRA —portada y título de cada lámina— y sans para lo que INFORMA —cuerpo,
% viñetas, índice de secciones y pie—. Ver \serifalab más abajo y su uso en
% \setbeamerfont{title/subtitle/frametitle/framesubtitle}.
%
% Elección de la familia, verificada y no adivinada. De las candidatas del
% encargo, kpsewhich confirmó instaladas: tgpagella, tgtermes, tgschola,
% tgbonum, mathpazo, charter (no estaban XCharter ni newtxtext). Se compilaron
% y renderizaron a PNG cuatro candidatas —Pagella, Schola, Charter, Bonum— con
% el título de lámina y la banda de portada a tamaño real de proyección
% (comparar.tex/.pdf, carpeta de trabajo temporal). Resultado:
%   · TeX Gyre Schola y TeX Gyre Bonum: demasiado pesadas y de baja apertura,
%     compiten con la banda de color en vez de flotar sobre ella.
%   · Charter: elegante pero de contraste bajo y aire muy contemporáneo; se
%     aleja del carácter clásico-académico del modelo.
%   · TeX Gyre Pagella (clon de Palatino): contraste moderado, ojo grande,
%     aperturas abiertas —el «a» y el «g» de dos pisos que se ven en el
%     modelo—, y legible a tamaño de título incluso hasta el fondo de un
%     salón. Es además, junto con TeX Gyre Heros, parte de la misma familia
%     TeX Gyre: ambas fueron diseñadas para convivir, con métricas y peso
%     visual pensados para mezclarse. GANADORA.
%
% Se invoca la familia NFSS "qpl" directamente con \fontfamily, SIN cargar
% \usepackage{tgpagella}: ese paquete redefine \rmdefault a qpl, y aunque acá
% no se usa \rmfamily/\textrm/\mathrm en ninguna lámina (verificado con
% grep), no hay razón para tocar un valor global que no nos pertenece.
% \fontfamily{qpl} basta: NFSS localiza sola el .fd de TeX Gyre Pagella la
% primera vez que se usa. \familydefault SIGUE siendo \sfdefault: el cuerpo
% del deck no cambia.
%
% BUG CORREGIDO — \serifalab llevaba SOLO \fontfamily{qpl}, sin \selectfont.
% \fontfamily por sí solo no cambia la fuente activa: solo prepara el
% registro interno de NFSS: el cambio real ocurre hasta el siguiente
% \selectfont. En \setbeamerfont{title/subtitle/frametitle/framesubtitle}
% esto no se notaba porque la clave "family=" de beamer llama a \selectfont
% ella sola, al final de \usebeamerfont (ver beamerbasefont.sty). Pero en el
% pie de página \serifalab se usaba suelto, después de \usebeamerfont{footline}
% —que ya había hecho SU propio \selectfont en sans— y como nunca se volvía a
% llamar \selectfont, el cambio de familia quedaba preparado pero nunca se
% aplicaba: el pie seguía viéndose en sans siempre, sin importar cuántas veces
% se recompilara. \serifalab ahora incluye su propio \selectfont, así que
% funciona igual se use suelto o dentro de family=.
\newcommand{\serifalab}{\fontfamily{qpl}\selectfont}

% ── MATEMÁTICAS ───────────────────────────────────────────────────────────────
% Cargadas por si una lámina necesita un subíndice o una fracción.
% Recordatorio del §7: cero fórmulas con sumatorias en las láminas.
\usepackage{amsmath}
\usepackage{amssymb}

% ── GRÁFICAS Y TABLAS ─────────────────────────────────────────────────────────
\usepackage{graphicx}
\usepackage{booktabs}
\usepackage{tabularx}
\usepackage{array}

% ── COLOR Y CAJAS ─────────────────────────────────────────────────────────────
\usepackage{xcolor}
% ══════════════════════════════════════════════════════════════════════════════
% PALETA DEL LABORATORIO DE R PARA CIENCIA POLÍTICA — ITAM
% Verde institucional #006847 sobre blanco
%
% ESTE ES EL ÚNICO LUGAR NORMATIVO DONDE VIVEN LOS HEXES.
% Si cambia un color, se cambia aquí y se propaga a mano a los otros dos
% archivos que lo heredan:
%   · estilo/tema_lab.R     → theme_lab() y las escalas de ggplot2
%   · estilo/estilo.scss    → el sitio Quarto
% Los tres tienen que decir lo mismo. Si no, las láminas, los gráficos y el
% sitio dejan de combinar y se nota.
% ══════════════════════════════════════════════════════════════════════════════

\definecolor{verdeITAM}{HTML}{006847}     % acento principal — verde institucional
\definecolor{verdeProfundo}{HTML}{00402C} % títulos y texto enfático
\definecolor{verdeClaro}{HTML}{A8D5BE}    % fondos suaves, resaltados
\definecolor{grisTexto}{HTML}{3D3D3D}     % cuerpo de texto
\definecolor{guinda}{HTML}{7B1113}        % acento secundario, uso escaso

% ── Derivados ────────────────────────────────────────────────────────────────
% No son colores nuevos: son diluciones de los anteriores. Se definen aquí para
% que nadie escriba "verdeClaro!15" a mano en una lámina y el curso termine con
% quince tonos distintos repartidos por ahí.

\colorlet{verdeVelo}{verdeClaro!22!white}   % fondo de bloque y de código
\colorlet{verdeTenue}{verdeClaro!55!white}  % resaltado dentro de un bloque
\colorlet{verdeSombra}{white!30!verdeITAM}  % texto atenuado SOBRE fondo verde
\colorlet{grisVelo}{grisTexto!7!white}      % fondos muy tenues, reglas
\colorlet{grisTenue}{grisTexto!45!white}    % texto terciario, pies
   % ← única fuente de los hexes
\usepackage{tcolorbox}
\tcbuselibrary{skins, breakable}

\usepackage{tikz}
\usetikzlibrary{calc}

% ── CÓDIGO EN LAS LÁMINAS ─────────────────────────────────────────────────────
% El código va en las láminas porque muchas personas no alcanzan a copiar del
% proyector a su pantalla. Se usa listings y no minted: minted exige Python,
% Pygments y compilar con -shell-escape, y eso es fricción que nadie debería
% tener que reproducir para compilar el material de un curso.
\usepackage{listings}

\lstdefinestyle{Rlab}{
  language         = R,
  % \small y no \scriptsize: esto se proyecta en un salón y hay gente hasta el
  % fondo. Si un bloque no cabe a \small, el bloque es demasiado largo para una
  % lámina — se parte en dos, no se encoge la letra.
  basicstyle       = \ttfamily\small\color{grisTexto},
  keywordstyle     = \color{verdeITAM}\bfseries,
  % Los comentarios NO van en gris tenue: en este curso el comentario es el
  % contenido, no una nota al margen.
  commentstyle     = \color{grisTexto}\itshape,
  stringstyle      = \color{guinda},
  identifierstyle  = \color{grisTexto},
  backgroundcolor  = \color{verdeVelo},
  frame            = leftline,
  framesep         = 6pt,
  rulecolor        = \color{verdeITAM},
  framerule        = 1.2pt,
  xleftmargin      = 8pt,
  showstringspaces = false,
  breaklines       = true,
  breakatwhitespace= true,
  columns          = fullflexible,
  keepspaces       = true,
  upquote          = true,
  aboveskip        = 0.7em,
  belowskip        = 0.7em,
  literate =
    {á}{{\'a}}1 {é}{{\'e}}1 {í}{{\'i}}1 {ó}{{\'o}}1 {ú}{{\'u}}1
    {Á}{{\'A}}1 {É}{{\'E}}1 {Í}{{\'I}}1 {Ó}{{\'O}}1 {Ú}{{\'U}}1
    {ñ}{{\~n}}1 {Ñ}{{\~N}}1 {¿}{{?`}}1 {¡}{{!`}}1
    % El pipe nativo no es palabra reservada de R para listings y no puede
    % declararse como keyword porque no es alfanumérico. Se colorea a mano.
    % Es el operador que más aparece en todo el curso: tiene que verse.
    {|>}{{\textcolor{verdeITAM}{\bfseries|>}}}2
}
\lstset{style=Rlab}

% Verbos del tidyverse. No son palabras reservadas de R, pero en este curso son
% el vocabulario central: si se ven siempre del mismo color, el grupo empieza a
% reconocerlos antes de saber qué hacen.
\lstset{morekeywords={tibble, glimpse, mutate, filter, select, arrange,
  summarise, group_by, left_join, anti_join, count, ggplot, aes, labs,
  geom_col, geom_point, geom_sf, geom_histogram, st_read, st_crs, here, feols,
  pivot_longer, pivot_wider, read_csv, read_rds, clean_names, slice_max,
  case_when, theme_lab, ggsave, modelsummary, facet_wrap}}

% Código en línea dentro de un párrafo: \cod{filter()}
\newcommand{\cod}[1]{{\ttfamily\color{verdeProfundo}#1}}

% Énfasis. NUNCA \textbf en el cuerpo: el énfasis va por color.
\newcommand{\ac}[1]{\textcolor{verdeITAM}{#1}}   % acento primario
\newcommand{\se}[1]{\textcolor{guinda}{#1}}      % acento secundario, uso escaso

% ── LAYOUT Y UTILIDADES ───────────────────────────────────────────────────────
\usepackage{multicol}
\usepackage{appendixnumberbeamer}

% ── HIPERENLACES — siempre al final ───────────────────────────────────────────
\usepackage{url}
\usepackage{hyperref}

% ══════════════════════════════════════════════════════════════════════════════
% TEMA
% ══════════════════════════════════════════════════════════════════════════════
\usetheme{default}
\usecolortheme{default}
\usefonttheme{professionalfonts}   % NO {serif}: el deck es sans en todo
\setbeamertemplate{navigation symbols}{}

% ── Colores ───────────────────────────────────────────────────────────────────
\setbeamercolor{normal text}{fg=grisTexto, bg=white}
\setbeamercolor{structure}{fg=verdeITAM}

\setbeamercolor{frametitle}{fg=verdeITAM, bg=white}
\setbeamercolor{framesubtitle}{fg=grisTenue}

\setbeamercolor{item}{fg=verdeITAM}
\setbeamercolor{subitem}{fg=verdeITAM}
\setbeamercolor{subsubitem}{fg=grisTenue}

\setbeamercolor{block title}{fg=white, bg=verdeITAM}
\setbeamercolor{block body}{fg=grisTexto, bg=verdeVelo}
\setbeamercolor{block title alerted}{fg=white, bg=guinda}
\setbeamercolor{block body alerted}{fg=grisTexto, bg=grisVelo}
\setbeamercolor{block title example}{fg=verdeProfundo, bg=verdeTenue}
\setbeamercolor{block body example}{fg=grisTexto, bg=verdeVelo}

% Índice de las láminas de sección: la actual en blanco, las demás atenuadas.
% Se fijan SOLO los dos colores y se deja la plantilla de beamer intacta.
% Redefinir \setbeamertemplate{section in toc} a mano hace que la lámina de
% sección salga en blanco: probado, y no vale la pena averiguar por qué.
\setbeamercolor{section in toc}{fg=white}
\setbeamercolor{section in toc shaded}{fg=verdeSombra}

% ── Tamaños ───────────────────────────────────────────────────────────────────
% title, subtitle, frametitle y framesubtitle llevan family=\serifalab: son
% los cuatro elementos que "nombran" (portada y título de lámina). Todo lo
% demás —author, institute, date, block body, footline, section in toc— se
% deja SIN family, así que hereda \familydefault, que sigue siendo sans.
\setbeamerfont{title}{family=\serifalab, size=\LARGE}
\setbeamerfont{subtitle}{family=\serifalab, size=\large}
\setbeamerfont{author}{size=\normalsize}
\setbeamerfont{institute}{size=\small}
\setbeamerfont{date}{size=\small}
\setbeamerfont{frametitle}{family=\serifalab, size=\Large}
\setbeamerfont{framesubtitle}{family=\serifalab, size=\small}
\setbeamerfont{block title}{size=\normalsize, series=\bfseries}
\setbeamerfont{block body}{size=\small}
\setbeamerfont{footline}{size=\tiny}
\setbeamerfont{section in toc}{size=\large}

% ── Viñetas ───────────────────────────────────────────────────────────────────
% Punto medio pequeño, como en los decks de seminario. Nada de triángulos ni
% de cuadrados: la viñeta no debe competir con el texto.
\setbeamertemplate{itemize item}{\raisebox{0.12ex}{\scriptsize\textcolor{verdeITAM}{$\bullet$}}}
\setbeamertemplate{itemize subitem}{\raisebox{0.12ex}{\tiny\textcolor{verdeITAM}{$\bullet$}}}
\setbeamertemplate{enumerate item}{\textcolor{verdeITAM}{\insertenumlabel.}}
\setbeamertemplate{caption}{\insertcaption}

% ══════════════════════════════════════════════════════════════════════════════
% ENCABEZADO — banda superior con el logo del ITAM sobre verde
% ══════════════════════════════════════════════════════════════════════════════
\newcommand{\rutaLogo}{../../estilo/logo_itam.png}

\setbeamercolor{cajalogo}{fg=white, bg=verdeITAM}
\setbeamercolor{bandacurso}{fg=grisTenue, bg=white}

% Se arma con beamercolorbox y NO con tikz. Un `\fill` dentro de una
% tikzpicture con `overlay` en la plantilla de headline no se pinta —el nodo
% de la imagen sí aparece, pero el rectángulo de fondo no—, y como el
% lettering del ITAM es blanco sobre transparente, el logo quedaba invisible
% sobre el blanco de la lámina. Verificado con pdftoppm a 300 ppp.
% Historial: 2.4ex -> 3.0ex -> 3.5ex. Cada vez crecí la caja casi en la misma
% proporción que la imagen para evitar que se desbordara, y el resultado fue
% que la fracción de la caja verde que ocupa el logo casi no cambiaba (pasó de
% 43% a 46% a 49%): un cambio real pero demasiado sutil para notarse a
% simple vista. Esta vez el salto es deliberadamente grande y la caja crece
% MENOS que proporcionalmente, para que el logo se vea inconfundiblemente más
% grande y no solo "técnicamente" más grande: sube a 5.0ex (+43% sobre 3.5ex),
% con ht=5.2ex/dp=2.8ex (antes 4.6/2.6). El logo pasa a ocupar ~62% de la caja
% (antes ~49%). Sigue en el mismo lugar: esquina superior derecha, centrado.
% Recalculo r = (ht - dp)/2 - altura_imagen/2 = (5.2-2.8)/2 - 5.0/2 = -1.3ex.
% La caja izquierda lleva EXACTAMENTE los mismos ht y dp, o las dos mitades
% del encabezado quedan a distinta altura.
\setbeamertemplate{headline}{%
  \leavevmode%
  \hbox{%
    \begin{beamercolorbox}[wd=0.815\paperwidth, ht=5.2ex, dp=2.8ex, leftskip=1.4em]{bandacurso}%
      {\tiny LABORATORIO DE R PARA CIENCIA POLÍTICA}%
    \end{beamercolorbox}%
    % El logo se dimensiona en ex y no en fracción del ancho de página: así
    % cabe siempre dentro de la caja (ht 5.2ex + dp 2.8ex = 8.0ex de alto) y
    % no se desborda por arriba. El desplazamiento lo centra verticalmente.
    \begin{beamercolorbox}[wd=0.185\paperwidth, ht=5.2ex, dp=2.8ex, center]{cajalogo}%
      \raisebox{-1.3ex}{\includegraphics[height=5.0ex]{\rutaLogo}}%
    \end{beamercolorbox}%
  }%
  {\color{grisTexto!22}\hrule height 0.4pt}%
  \vskip0.7em%
}

% ══════════════════════════════════════════════════════════════════════════════
% PIE — sección a la izquierda, iniciales y folio a la derecha
% ══════════════════════════════════════════════════════════════════════════════
% Iniciales y folio ~1mm más a la izquierda (Cambio 3).
%
% DIAGNÓSTICO — antes de tocar el valor, se descubrió que "rightskip" en esta
% caja NUNCA había tenido efecto, ni con 1.4em ni con ningún otro valor: la
% clave "right" de beamercolorbox (beamerbasecolor.sty) hace
% \def\beamer@colbox@rs{0pt} y, como en la lista de opciones "right" iba
% DESPUÉS de "rightskip=1.4em", pisaba ese valor y lo dejaba en 0pt fijo.
% Medido en el PDF ya compilado: el «0» de «1/20» quedaba a 0.25mm del borde
% derecho de la lámina, no al 1.4em (≈5mm) que el código sugería. Es un bug
% latente de origen, no algo que haya introducido este cambio.
%
% ARREGLO — se invierte el orden ("right" primero, "rightskip" después) para
% que nuestro valor sea el que sobrevive. Y como el punto de partida real era
% 0pt (no 1.4em), para lograr el desplazamiento sutil de ~1mm que pidió EMG
% —"milimétricamente", no un recorte grande— el valor correcto NO es 2.2em
% (eso habría movido el bloque ~7.8mm, ocho veces más de lo pedido): es
% 0.45em, calibrado por render y medición en píxeles (antes: "0" de "1/20" a
% 0.25mm del borde; con 0.45em: a 1.02mm — ver reporte). La caja izquierda
% (sección) no se toca.
% Iniciales y folio («EMG | N/total») en serifa: \serifalab se aplica SOLO
% dentro de esta caja, después de \usebeamerfont{footline} (que solo trae el
% tamaño \tiny, sin family). El nombre de la sección de la caja izquierda
% NO lleva \serifalab y se queda en sans, como el resto del pie.
\setbeamertemplate{footline}{%
  \hbox{%
    \begin{beamercolorbox}[wd=0.62\paperwidth, ht=2.6ex, dp=1.4ex, leftskip=1.4em]{}%
      {\usebeamerfont{footline}\color{verdeITAM}\insertsectionhead}%
    \end{beamercolorbox}%
    \begin{beamercolorbox}[wd=0.38\paperwidth, ht=2.6ex, dp=1.4ex, right, rightskip=0.45em]{}%
      {\usebeamerfont{footline}\serifalab\color{grisTenue}%
       EMG\hspace{0.9em}\textbar\hspace{0.9em}%
       \insertframenumber{}\,/\,\inserttotalframenumber}%
    \end{beamercolorbox}%
  }%
  \vskip2pt%
}

% ══════════════════════════════════════════════════════════════════════════════
% TÍTULO DE FRAME — en verde, alineado a la izquierda, SIN regla debajo
% ══════════════════════════════════════════════════════════════════════════════
\setbeamertemplate{frametitle}{%
  \vspace{0.6em}
  \hspace{-0.4em}%
  \begin{minipage}{0.95\textwidth}
    \raggedright
    {\usebeamerfont{frametitle}\color{verdeITAM}\insertframetitle\par}%
    \ifx\insertframesubtitle\@empty\else%
      \vspace{0.15em}%
      {\usebeamerfont{framesubtitle}\color{grisTenue}\insertframesubtitle\par}%
    \fi
  \end{minipage}
  \vspace{0.5em}
}

% ══════════════════════════════════════════════════════════════════════════════
% PORTADA — banda verde con el título, datos debajo sobre blanco
% Título y subtítulo salen en serifa porque \usebeamerfont{title} y
% \usebeamerfont{subtitle} ya llevan family=\serifalab (ver §Tamaños): no hace
% falta tocar \inserttitle/\insertsubtitle aquí abajo. Autoría, adscripción y
% fecha usan \usebeamerfont{author/institute/date}, que NO llevan family y por
% lo tanto heredan \familydefault (sans) sin cambios.
% ══════════════════════════════════════════════════════════════════════════════
\setbeamertemplate{title page}{%
  \begin{tikzpicture}[remember picture, overlay]
    % Banda de color con el título dentro
    \fill[verdeITAM]
      ([yshift=-0.145\paperheight]current page.north west)
      rectangle ([yshift=-0.455\paperheight]current page.north east);
    \node[anchor=center, text width=0.86\paperwidth, align=center, inner sep=0pt]
      at ([yshift=-0.265\paperheight]current page.north)
      {\usebeamerfont{title}\color{white}\inserttitle};
    \node[anchor=center, text width=0.80\paperwidth, align=center, inner sep=0pt]
      at ([yshift=-0.385\paperheight]current page.north)
      {\usebeamerfont{subtitle}\color{verdeClaro}\insertsubtitle};

    % Autoría, adscripción y fecha
    \node[anchor=north, text width=0.86\paperwidth, align=center, inner sep=0pt]
      at ([yshift=-0.545\paperheight]current page.north)
      {\usebeamerfont{author}\color{grisTexto}\insertauthor};
    \node[anchor=north, text width=0.86\paperwidth, align=center, inner sep=0pt]
      at ([yshift=-0.625\paperheight]current page.north)
      {\usebeamerfont{institute}\color{grisTexto}\insertinstitute};
    \node[anchor=north, text width=0.86\paperwidth, align=center, inner sep=0pt]
      at ([yshift=-0.790\paperheight]current page.north)
      {\usebeamerfont{date}\color{grisTenue}\insertdate};
  \end{tikzpicture}%
}

% ══════════════════════════════════════════════════════════════════════════════
% LÁMINA DE SECCIÓN — fondo verde completo, índice con la sección actual en
% blanco y las demás atenuadas. Se dispara sola en cada \section{}.
% ══════════════════════════════════════════════════════════════════════════════
\AtBeginSection[]{%
  {%
  \setbeamercolor{background canvas}{bg=verdeITAM}
  % Las entradas del índice son hipervínculos, y con colorlinks=true hyperref
  % las pinta de linkcolor —que es el mismo verde del fondo—, así que la
  % sección actual desaparecía. Con hidelinks, hyperref deja de colorear y de
  % dibujar marcos, y entonces sí mandan los colores de beamer: la sección
  % actual en blanco y las demás atenuadas.
  % Diagnosticado a fuerza de bisecar el preámbulo; no es nada obvio.
  \hypersetup{hidelinks}
  \setbeamertemplate{headline}{}
  % Mismo arreglo y mismo valor que el footline normal (Cambio 3, ver
  % diagnóstico completo ahí arriba): "right" antes de "rightskip=0.45em", o la
  % clave "right" pisa nuestro valor y lo deja en 0pt. Si aquí se deja
  % distinto del footline normal, el pie de las láminas de sección queda
  % desalineado con el del resto.
  % Mismo \serifalab que el footline normal, y por la misma razón: solo en
  % «EMG | N/total», nunca en el nombre de la sección.
  \setbeamertemplate{footline}{%
    \hbox{%
      \begin{beamercolorbox}[wd=0.62\paperwidth, ht=2.6ex, dp=1.4ex, leftskip=1.4em]{}%
        {\usebeamerfont{footline}\color{white}\insertsectionhead}%
      \end{beamercolorbox}%
      \begin{beamercolorbox}[wd=0.38\paperwidth, ht=2.6ex, dp=1.4ex, right, rightskip=0.45em]{}%
        {\usebeamerfont{footline}\serifalab\color{verdeSombra}%
         EMG\hspace{0.9em}\textbar\hspace{0.9em}%
         \insertframenumber{}\,/\,\inserttotalframenumber}%
      \end{beamercolorbox}%
    }%
    \vskip2pt%
  }
  \begin{frame}[plain, noframenumbering]
    \vspace{1.4em}
    \begin{minipage}{0.86\textwidth}
      \hspace{1.2em}%
      \begin{minipage}{0.95\textwidth}
        \tableofcontents[currentsection, sectionstyle=show/shaded, hideallsubsections]
      \end{minipage}
    \end{minipage}
  \end{frame}
  }%
}

% ══════════════════════════════════════════════════════════════════════════════
% ENTORNOS PROPIOS DEL CURSO
% ══════════════════════════════════════════════════════════════════════════════

% Rótulo en versalitas falsas — tgheros sí tiene versalitas, pero el rótulo se
% ve mejor en mayúsculas con la letra pequeña.
\newcommand{\etiqueta}[1]{{\footnotesize\color{grisTenue}\MakeUppercase{#1}}}

% ── ¿Por qué nos importa esto? ────────────────────────────────────────────────
% Abre todas las sesiones: plantea la pregunta politológica antes de que
% aparezca una sola función.
\newcommand{\porquenosimporta}[2]{%
  \begin{frame}[plain, noframenumbering]
    \vfill
    \begin{center}
      \etiqueta{¿Por qué nos importa esto?}\\[1.3em]
      {\LARGE\color{verdeProfundo} #1}\\[1.5em]
      \begin{minipage}{0.80\textwidth}
        \centering\normalsize\color{grisTexto} #2
      \end{minipage}
    \end{center}
    \vfill
  \end{frame}%
}

% ── Lo que ya sabemos hacer ───────────────────────────────────────────────────
% Cierra todas las sesiones. Exactamente tres viñetas: es la recompensa.
\newcommand{\yasabemos}[3]{%
  \begin{frame}{Lo que ya sabemos hacer}
    \vspace{1.0em}
    \begin{itemize}\setlength{\itemsep}{1.4em}
      \item #1
      \item #2
      \item #3
    \end{itemize}
    \vfill
  \end{frame}%
}

% ── Adelanto de la sesión siguiente ───────────────────────────────────────────
\newcommand{\lasiguiente}[2]{%
  \begin{frame}[plain, noframenumbering]
    \vfill
    \begin{center}
      \etiqueta{La semana que viene}\\[1.3em]
      {\LARGE\color{verdeITAM} #1}\\[1.5em]
      \begin{minipage}{0.80\textwidth}
        \centering\normalsize\color{grisTexto} #2
      \end{minipage}
    \end{center}
    \vfill
  \end{frame}%
}

% ── Caja de error anticipado ──────────────────────────────────────────────────
% Los errores son parte del temario. Esta caja los presenta ANTES de que
% ocurran, con su traducción del inglés y su arreglo.
% "no shadow" y "sharp corners" son deliberados: nada de sombras ni degradados.
\newtcolorbox{errorbox}[1]{
  enhanced, breakable, no shadow, sharp corners,
  colback   = grisVelo,
  colframe  = guinda,
  boxrule   = 0pt, leftrule = 2.5pt,
  toprule = 0pt, bottomrule = 0pt, rightrule = 0pt,
  fontupper = \small\color{grisTexto},
  title     = {\color{guinda}\bfseries\small\ttfamily #1},
  colbacktitle = grisVelo,
  coltitle  = guinda,
  titlerule = 0pt
}

% ── Caja del salto politólogo ─────────────────────────────────────────────────
% El §7 pide que cada sesión contenga al menos un momento donde la intuición
% política se convierta en pregunta empírica, y que ese momento se nombre.
% Esta caja es el lugar donde se nombra.
\newtcolorbox{saltobox}{
  enhanced, breakable, no shadow, sharp corners,
  colback   = verdeVelo,
  colframe  = verdeITAM,
  boxrule   = 0pt, leftrule = 2.5pt,
  toprule = 0pt, bottomrule = 0pt, rightrule = 0pt,
  fontupper = \small\color{grisTexto},
  title     = {\color{verdeProfundo}\bfseries\small El salto},
  colbacktitle = verdeVelo,
  coltitle  = verdeProfundo,
  titlerule = 0pt
}

% ══════════════════════════════════════════════════════════════════════════════
% METADATOS FIJOS
% ══════════════════════════════════════════════════════════════════════════════
\author[E. Miranda]{Emiliano Miranda González}
\institute[ITAM]{%
  Instituto Tecnológico Autónomo de México\\
  Departamento Académico de Ciencia Política%
}

\hypersetup{
  colorlinks  = true,
  linkcolor   = verdeITAM,
  urlcolor    = verdeITAM,
  citecolor   = grisTenue,
  pdfauthor   = {Emiliano Miranda González},
  pdfsubject  = {Laboratorio de R para Ciencia Política — ITAM},
  pdfkeywords = {R, ciencia política, ITAM, métodos cuantitativos},
  unicode     = true
}

%
% Se compila con pdflatex, DOS VECES (los nodos de tikz con `remember picture`
% necesitan la segunda pasada para colocarse bien).
% ══════════════════════════════════════════════════════════════════════════════

\documentclass[
  aspectratio=43,
  t,
  10pt
]{beamer}

% ── TIPOGRAFÍA ────────────────────────────────────────────────────────────────
% El cuerpo del deck sigue siendo sans en su totalidad. TeX Gyre Heros es un
% clon de Helvetica en Type 1 con cobertura T1 completa —incluidas cursivas y
% versalitas reales—, presente en cualquier distribución de TeX Live. lmodern
% va ANTES para quedarse con la monoespaciada: sin él, los comentarios en
% cursiva de los bloques de código salen como fuente de mapa de bits y se ven
% borrosos al proyectar.
%
% Si en la máquina de quien compila están instaladas Fira Sans o Lato —que son
% las humanistas que suelen usar los decks de seminario de ciencia política—,
% basta sustituir la línea de tgheros por \usepackage[sfdefault]{FiraSans} o
% \usepackage[default]{lato}. El resto del preámbulo no cambia.
\usepackage[T1]{fontenc}
\usepackage[utf8]{inputenc}
\usepackage{lmodern}
\usepackage{tgheros}
\renewcommand{\familydefault}{\sfdefault}
\usepackage{microtype}

% ── SERIFA PARA TÍTULOS (mezcla del deck de seminario) ─────────────────────────
% EMG pidió la mezcla de los decks de seminario de ciencia política
% internacional (ref.: deck de Gary King, IQSS/Harvard): serifa para lo que
% NOMBRA —portada y título de cada lámina— y sans para lo que INFORMA —cuerpo,
% viñetas, índice de secciones y pie—. Ver \serifalab más abajo y su uso en
% \setbeamerfont{title/subtitle/frametitle/framesubtitle}.
%
% Elección de la familia, verificada y no adivinada. De las candidatas del
% encargo, kpsewhich confirmó instaladas: tgpagella, tgtermes, tgschola,
% tgbonum, mathpazo, charter (no estaban XCharter ni newtxtext). Se compilaron
% y renderizaron a PNG cuatro candidatas —Pagella, Schola, Charter, Bonum— con
% el título de lámina y la banda de portada a tamaño real de proyección
% (comparar.tex/.pdf, carpeta de trabajo temporal). Resultado:
%   · TeX Gyre Schola y TeX Gyre Bonum: demasiado pesadas y de baja apertura,
%     compiten con la banda de color en vez de flotar sobre ella.
%   · Charter: elegante pero de contraste bajo y aire muy contemporáneo; se
%     aleja del carácter clásico-académico del modelo.
%   · TeX Gyre Pagella (clon de Palatino): contraste moderado, ojo grande,
%     aperturas abiertas —el «a» y el «g» de dos pisos que se ven en el
%     modelo—, y legible a tamaño de título incluso hasta el fondo de un
%     salón. Es además, junto con TeX Gyre Heros, parte de la misma familia
%     TeX Gyre: ambas fueron diseñadas para convivir, con métricas y peso
%     visual pensados para mezclarse. GANADORA.
%
% Se invoca la familia NFSS "qpl" directamente con \fontfamily, SIN cargar
% \usepackage{tgpagella}: ese paquete redefine \rmdefault a qpl, y aunque acá
% no se usa \rmfamily/\textrm/\mathrm en ninguna lámina (verificado con
% grep), no hay razón para tocar un valor global que no nos pertenece.
% \fontfamily{qpl} basta: NFSS localiza sola el .fd de TeX Gyre Pagella la
% primera vez que se usa. \familydefault SIGUE siendo \sfdefault: el cuerpo
% del deck no cambia.
%
% BUG CORREGIDO — \serifalab llevaba SOLO \fontfamily{qpl}, sin \selectfont.
% \fontfamily por sí solo no cambia la fuente activa: solo prepara el
% registro interno de NFSS: el cambio real ocurre hasta el siguiente
% \selectfont. En \setbeamerfont{title/subtitle/frametitle/framesubtitle}
% esto no se notaba porque la clave "family=" de beamer llama a \selectfont
% ella sola, al final de \usebeamerfont (ver beamerbasefont.sty). Pero en el
% pie de página \serifalab se usaba suelto, después de \usebeamerfont{footline}
% —que ya había hecho SU propio \selectfont en sans— y como nunca se volvía a
% llamar \selectfont, el cambio de familia quedaba preparado pero nunca se
% aplicaba: el pie seguía viéndose en sans siempre, sin importar cuántas veces
% se recompilara. \serifalab ahora incluye su propio \selectfont, así que
% funciona igual se use suelto o dentro de family=.
\newcommand{\serifalab}{\fontfamily{qpl}\selectfont}

% ── MATEMÁTICAS ───────────────────────────────────────────────────────────────
% Cargadas por si una lámina necesita un subíndice o una fracción.
% Recordatorio del §7: cero fórmulas con sumatorias en las láminas.
\usepackage{amsmath}
\usepackage{amssymb}

% ── GRÁFICAS Y TABLAS ─────────────────────────────────────────────────────────
\usepackage{graphicx}
\usepackage{booktabs}
\usepackage{tabularx}
\usepackage{array}

% ── COLOR Y CAJAS ─────────────────────────────────────────────────────────────
\usepackage{xcolor}
% ══════════════════════════════════════════════════════════════════════════════
% PALETA DEL LABORATORIO DE R PARA CIENCIA POLÍTICA — ITAM
% Verde institucional #006847 sobre blanco
%
% ESTE ES EL ÚNICO LUGAR NORMATIVO DONDE VIVEN LOS HEXES.
% Si cambia un color, se cambia aquí y se propaga a mano a los otros dos
% archivos que lo heredan:
%   · estilo/tema_lab.R     → theme_lab() y las escalas de ggplot2
%   · estilo/estilo.scss    → el sitio Quarto
% Los tres tienen que decir lo mismo. Si no, las láminas, los gráficos y el
% sitio dejan de combinar y se nota.
% ══════════════════════════════════════════════════════════════════════════════

\definecolor{verdeITAM}{HTML}{006847}     % acento principal — verde institucional
\definecolor{verdeProfundo}{HTML}{00402C} % títulos y texto enfático
\definecolor{verdeClaro}{HTML}{A8D5BE}    % fondos suaves, resaltados
\definecolor{grisTexto}{HTML}{3D3D3D}     % cuerpo de texto
\definecolor{guinda}{HTML}{7B1113}        % acento secundario, uso escaso

% ── Derivados ────────────────────────────────────────────────────────────────
% No son colores nuevos: son diluciones de los anteriores. Se definen aquí para
% que nadie escriba "verdeClaro!15" a mano en una lámina y el curso termine con
% quince tonos distintos repartidos por ahí.

\colorlet{verdeVelo}{verdeClaro!22!white}   % fondo de bloque y de código
\colorlet{verdeTenue}{verdeClaro!55!white}  % resaltado dentro de un bloque
\colorlet{verdeSombra}{white!30!verdeITAM}  % texto atenuado SOBRE fondo verde
\colorlet{grisVelo}{grisTexto!7!white}      % fondos muy tenues, reglas
\colorlet{grisTenue}{grisTexto!45!white}    % texto terciario, pies
   % ← única fuente de los hexes
\usepackage{tcolorbox}
\tcbuselibrary{skins, breakable}

\usepackage{tikz}
\usetikzlibrary{calc}

% ── CÓDIGO EN LAS LÁMINAS ─────────────────────────────────────────────────────
% El código va en las láminas porque muchas personas no alcanzan a copiar del
% proyector a su pantalla. Se usa listings y no minted: minted exige Python,
% Pygments y compilar con -shell-escape, y eso es fricción que nadie debería
% tener que reproducir para compilar el material de un curso.
\usepackage{listings}

\lstdefinestyle{Rlab}{
  language         = R,
  % \small y no \scriptsize: esto se proyecta en un salón y hay gente hasta el
  % fondo. Si un bloque no cabe a \small, el bloque es demasiado largo para una
  % lámina — se parte en dos, no se encoge la letra.
  basicstyle       = \ttfamily\small\color{grisTexto},
  keywordstyle     = \color{verdeITAM}\bfseries,
  % Los comentarios NO van en gris tenue: en este curso el comentario es el
  % contenido, no una nota al margen.
  commentstyle     = \color{grisTexto}\itshape,
  stringstyle      = \color{guinda},
  identifierstyle  = \color{grisTexto},
  backgroundcolor  = \color{verdeVelo},
  frame            = leftline,
  framesep         = 6pt,
  rulecolor        = \color{verdeITAM},
  framerule        = 1.2pt,
  xleftmargin      = 8pt,
  showstringspaces = false,
  breaklines       = true,
  breakatwhitespace= true,
  columns          = fullflexible,
  keepspaces       = true,
  upquote          = true,
  aboveskip        = 0.7em,
  belowskip        = 0.7em,
  literate =
    {á}{{\'a}}1 {é}{{\'e}}1 {í}{{\'i}}1 {ó}{{\'o}}1 {ú}{{\'u}}1
    {Á}{{\'A}}1 {É}{{\'E}}1 {Í}{{\'I}}1 {Ó}{{\'O}}1 {Ú}{{\'U}}1
    {ñ}{{\~n}}1 {Ñ}{{\~N}}1 {¿}{{?`}}1 {¡}{{!`}}1
    % El pipe nativo no es palabra reservada de R para listings y no puede
    % declararse como keyword porque no es alfanumérico. Se colorea a mano.
    % Es el operador que más aparece en todo el curso: tiene que verse.
    {|>}{{\textcolor{verdeITAM}{\bfseries|>}}}2
}
\lstset{style=Rlab}

% Verbos del tidyverse. No son palabras reservadas de R, pero en este curso son
% el vocabulario central: si se ven siempre del mismo color, el grupo empieza a
% reconocerlos antes de saber qué hacen.
\lstset{morekeywords={tibble, glimpse, mutate, filter, select, arrange,
  summarise, group_by, left_join, anti_join, count, ggplot, aes, labs,
  geom_col, geom_point, geom_sf, geom_histogram, st_read, st_crs, here, feols,
  pivot_longer, pivot_wider, read_csv, read_rds, clean_names, slice_max,
  case_when, theme_lab, ggsave, modelsummary, facet_wrap}}

% Código en línea dentro de un párrafo: \cod{filter()}
\newcommand{\cod}[1]{{\ttfamily\color{verdeProfundo}#1}}

% Énfasis. NUNCA \textbf en el cuerpo: el énfasis va por color.
\newcommand{\ac}[1]{\textcolor{verdeITAM}{#1}}   % acento primario
\newcommand{\se}[1]{\textcolor{guinda}{#1}}      % acento secundario, uso escaso

% ── LAYOUT Y UTILIDADES ───────────────────────────────────────────────────────
\usepackage{multicol}
\usepackage{appendixnumberbeamer}

% ── HIPERENLACES — siempre al final ───────────────────────────────────────────
\usepackage{url}
\usepackage{hyperref}

% ══════════════════════════════════════════════════════════════════════════════
% TEMA
% ══════════════════════════════════════════════════════════════════════════════
\usetheme{default}
\usecolortheme{default}
\usefonttheme{professionalfonts}   % NO {serif}: el deck es sans en todo
\setbeamertemplate{navigation symbols}{}

% ── Colores ───────────────────────────────────────────────────────────────────
\setbeamercolor{normal text}{fg=grisTexto, bg=white}
\setbeamercolor{structure}{fg=verdeITAM}

\setbeamercolor{frametitle}{fg=verdeITAM, bg=white}
\setbeamercolor{framesubtitle}{fg=grisTenue}

\setbeamercolor{item}{fg=verdeITAM}
\setbeamercolor{subitem}{fg=verdeITAM}
\setbeamercolor{subsubitem}{fg=grisTenue}

\setbeamercolor{block title}{fg=white, bg=verdeITAM}
\setbeamercolor{block body}{fg=grisTexto, bg=verdeVelo}
\setbeamercolor{block title alerted}{fg=white, bg=guinda}
\setbeamercolor{block body alerted}{fg=grisTexto, bg=grisVelo}
\setbeamercolor{block title example}{fg=verdeProfundo, bg=verdeTenue}
\setbeamercolor{block body example}{fg=grisTexto, bg=verdeVelo}

% Índice de las láminas de sección: la actual en blanco, las demás atenuadas.
% Se fijan SOLO los dos colores y se deja la plantilla de beamer intacta.
% Redefinir \setbeamertemplate{section in toc} a mano hace que la lámina de
% sección salga en blanco: probado, y no vale la pena averiguar por qué.
\setbeamercolor{section in toc}{fg=white}
\setbeamercolor{section in toc shaded}{fg=verdeSombra}

% ── Tamaños ───────────────────────────────────────────────────────────────────
% title, subtitle, frametitle y framesubtitle llevan family=\serifalab: son
% los cuatro elementos que "nombran" (portada y título de lámina). Todo lo
% demás —author, institute, date, block body, footline, section in toc— se
% deja SIN family, así que hereda \familydefault, que sigue siendo sans.
\setbeamerfont{title}{family=\serifalab, size=\LARGE}
\setbeamerfont{subtitle}{family=\serifalab, size=\large}
\setbeamerfont{author}{size=\normalsize}
\setbeamerfont{institute}{size=\small}
\setbeamerfont{date}{size=\small}
\setbeamerfont{frametitle}{family=\serifalab, size=\Large}
\setbeamerfont{framesubtitle}{family=\serifalab, size=\small}
\setbeamerfont{block title}{size=\normalsize, series=\bfseries}
\setbeamerfont{block body}{size=\small}
\setbeamerfont{footline}{size=\tiny}
\setbeamerfont{section in toc}{size=\large}

% ── Viñetas ───────────────────────────────────────────────────────────────────
% Punto medio pequeño, como en los decks de seminario. Nada de triángulos ni
% de cuadrados: la viñeta no debe competir con el texto.
\setbeamertemplate{itemize item}{\raisebox{0.12ex}{\scriptsize\textcolor{verdeITAM}{$\bullet$}}}
\setbeamertemplate{itemize subitem}{\raisebox{0.12ex}{\tiny\textcolor{verdeITAM}{$\bullet$}}}
\setbeamertemplate{enumerate item}{\textcolor{verdeITAM}{\insertenumlabel.}}
\setbeamertemplate{caption}{\insertcaption}

% ══════════════════════════════════════════════════════════════════════════════
% ENCABEZADO — banda superior con el logo del ITAM sobre verde
% ══════════════════════════════════════════════════════════════════════════════
\newcommand{\rutaLogo}{../../estilo/logo_itam.png}

\setbeamercolor{cajalogo}{fg=white, bg=verdeITAM}
\setbeamercolor{bandacurso}{fg=grisTenue, bg=white}

% Se arma con beamercolorbox y NO con tikz. Un `\fill` dentro de una
% tikzpicture con `overlay` en la plantilla de headline no se pinta —el nodo
% de la imagen sí aparece, pero el rectángulo de fondo no—, y como el
% lettering del ITAM es blanco sobre transparente, el logo quedaba invisible
% sobre el blanco de la lámina. Verificado con pdftoppm a 300 ppp.
% Historial: 2.4ex -> 3.0ex -> 3.5ex. Cada vez crecí la caja casi en la misma
% proporción que la imagen para evitar que se desbordara, y el resultado fue
% que la fracción de la caja verde que ocupa el logo casi no cambiaba (pasó de
% 43% a 46% a 49%): un cambio real pero demasiado sutil para notarse a
% simple vista. Esta vez el salto es deliberadamente grande y la caja crece
% MENOS que proporcionalmente, para que el logo se vea inconfundiblemente más
% grande y no solo "técnicamente" más grande: sube a 5.0ex (+43% sobre 3.5ex),
% con ht=5.2ex/dp=2.8ex (antes 4.6/2.6). El logo pasa a ocupar ~62% de la caja
% (antes ~49%). Sigue en el mismo lugar: esquina superior derecha, centrado.
% Recalculo r = (ht - dp)/2 - altura_imagen/2 = (5.2-2.8)/2 - 5.0/2 = -1.3ex.
% La caja izquierda lleva EXACTAMENTE los mismos ht y dp, o las dos mitades
% del encabezado quedan a distinta altura.
\setbeamertemplate{headline}{%
  \leavevmode%
  \hbox{%
    \begin{beamercolorbox}[wd=0.815\paperwidth, ht=5.2ex, dp=2.8ex, leftskip=1.4em]{bandacurso}%
      {\tiny LABORATORIO DE R PARA CIENCIA POLÍTICA}%
    \end{beamercolorbox}%
    % El logo se dimensiona en ex y no en fracción del ancho de página: así
    % cabe siempre dentro de la caja (ht 5.2ex + dp 2.8ex = 8.0ex de alto) y
    % no se desborda por arriba. El desplazamiento lo centra verticalmente.
    \begin{beamercolorbox}[wd=0.185\paperwidth, ht=5.2ex, dp=2.8ex, center]{cajalogo}%
      \raisebox{-1.3ex}{\includegraphics[height=5.0ex]{\rutaLogo}}%
    \end{beamercolorbox}%
  }%
  {\color{grisTexto!22}\hrule height 0.4pt}%
  \vskip0.7em%
}

% ══════════════════════════════════════════════════════════════════════════════
% PIE — sección a la izquierda, iniciales y folio a la derecha
% ══════════════════════════════════════════════════════════════════════════════
% Iniciales y folio ~1mm más a la izquierda (Cambio 3).
%
% DIAGNÓSTICO — antes de tocar el valor, se descubrió que "rightskip" en esta
% caja NUNCA había tenido efecto, ni con 1.4em ni con ningún otro valor: la
% clave "right" de beamercolorbox (beamerbasecolor.sty) hace
% \def\beamer@colbox@rs{0pt} y, como en la lista de opciones "right" iba
% DESPUÉS de "rightskip=1.4em", pisaba ese valor y lo dejaba en 0pt fijo.
% Medido en el PDF ya compilado: el «0» de «1/20» quedaba a 0.25mm del borde
% derecho de la lámina, no al 1.4em (≈5mm) que el código sugería. Es un bug
% latente de origen, no algo que haya introducido este cambio.
%
% ARREGLO — se invierte el orden ("right" primero, "rightskip" después) para
% que nuestro valor sea el que sobrevive. Y como el punto de partida real era
% 0pt (no 1.4em), para lograr el desplazamiento sutil de ~1mm que pidió EMG
% —"milimétricamente", no un recorte grande— el valor correcto NO es 2.2em
% (eso habría movido el bloque ~7.8mm, ocho veces más de lo pedido): es
% 0.45em, calibrado por render y medición en píxeles (antes: "0" de "1/20" a
% 0.25mm del borde; con 0.45em: a 1.02mm — ver reporte). La caja izquierda
% (sección) no se toca.
% Iniciales y folio («EMG | N/total») en serifa: \serifalab se aplica SOLO
% dentro de esta caja, después de \usebeamerfont{footline} (que solo trae el
% tamaño \tiny, sin family). El nombre de la sección de la caja izquierda
% NO lleva \serifalab y se queda en sans, como el resto del pie.
\setbeamertemplate{footline}{%
  \hbox{%
    \begin{beamercolorbox}[wd=0.62\paperwidth, ht=2.6ex, dp=1.4ex, leftskip=1.4em]{}%
      {\usebeamerfont{footline}\color{verdeITAM}\insertsectionhead}%
    \end{beamercolorbox}%
    \begin{beamercolorbox}[wd=0.38\paperwidth, ht=2.6ex, dp=1.4ex, right, rightskip=0.45em]{}%
      {\usebeamerfont{footline}\serifalab\color{grisTenue}%
       EMG\hspace{0.9em}\textbar\hspace{0.9em}%
       \insertframenumber{}\,/\,\inserttotalframenumber}%
    \end{beamercolorbox}%
  }%
  \vskip2pt%
}

% ══════════════════════════════════════════════════════════════════════════════
% TÍTULO DE FRAME — en verde, alineado a la izquierda, SIN regla debajo
% ══════════════════════════════════════════════════════════════════════════════
\setbeamertemplate{frametitle}{%
  \vspace{0.6em}
  \hspace{-0.4em}%
  \begin{minipage}{0.95\textwidth}
    \raggedright
    {\usebeamerfont{frametitle}\color{verdeITAM}\insertframetitle\par}%
    \ifx\insertframesubtitle\@empty\else%
      \vspace{0.15em}%
      {\usebeamerfont{framesubtitle}\color{grisTenue}\insertframesubtitle\par}%
    \fi
  \end{minipage}
  \vspace{0.5em}
}

% ══════════════════════════════════════════════════════════════════════════════
% PORTADA — banda verde con el título, datos debajo sobre blanco
% Título y subtítulo salen en serifa porque \usebeamerfont{title} y
% \usebeamerfont{subtitle} ya llevan family=\serifalab (ver §Tamaños): no hace
% falta tocar \inserttitle/\insertsubtitle aquí abajo. Autoría, adscripción y
% fecha usan \usebeamerfont{author/institute/date}, que NO llevan family y por
% lo tanto heredan \familydefault (sans) sin cambios.
% ══════════════════════════════════════════════════════════════════════════════
\setbeamertemplate{title page}{%
  \begin{tikzpicture}[remember picture, overlay]
    % Banda de color con el título dentro
    \fill[verdeITAM]
      ([yshift=-0.145\paperheight]current page.north west)
      rectangle ([yshift=-0.455\paperheight]current page.north east);
    \node[anchor=center, text width=0.86\paperwidth, align=center, inner sep=0pt]
      at ([yshift=-0.265\paperheight]current page.north)
      {\usebeamerfont{title}\color{white}\inserttitle};
    \node[anchor=center, text width=0.80\paperwidth, align=center, inner sep=0pt]
      at ([yshift=-0.385\paperheight]current page.north)
      {\usebeamerfont{subtitle}\color{verdeClaro}\insertsubtitle};

    % Autoría, adscripción y fecha
    \node[anchor=north, text width=0.86\paperwidth, align=center, inner sep=0pt]
      at ([yshift=-0.545\paperheight]current page.north)
      {\usebeamerfont{author}\color{grisTexto}\insertauthor};
    \node[anchor=north, text width=0.86\paperwidth, align=center, inner sep=0pt]
      at ([yshift=-0.625\paperheight]current page.north)
      {\usebeamerfont{institute}\color{grisTexto}\insertinstitute};
    \node[anchor=north, text width=0.86\paperwidth, align=center, inner sep=0pt]
      at ([yshift=-0.790\paperheight]current page.north)
      {\usebeamerfont{date}\color{grisTenue}\insertdate};
  \end{tikzpicture}%
}

% ══════════════════════════════════════════════════════════════════════════════
% LÁMINA DE SECCIÓN — fondo verde completo, índice con la sección actual en
% blanco y las demás atenuadas. Se dispara sola en cada \section{}.
% ══════════════════════════════════════════════════════════════════════════════
\AtBeginSection[]{%
  {%
  \setbeamercolor{background canvas}{bg=verdeITAM}
  % Las entradas del índice son hipervínculos, y con colorlinks=true hyperref
  % las pinta de linkcolor —que es el mismo verde del fondo—, así que la
  % sección actual desaparecía. Con hidelinks, hyperref deja de colorear y de
  % dibujar marcos, y entonces sí mandan los colores de beamer: la sección
  % actual en blanco y las demás atenuadas.
  % Diagnosticado a fuerza de bisecar el preámbulo; no es nada obvio.
  \hypersetup{hidelinks}
  \setbeamertemplate{headline}{}
  % Mismo arreglo y mismo valor que el footline normal (Cambio 3, ver
  % diagnóstico completo ahí arriba): "right" antes de "rightskip=0.45em", o la
  % clave "right" pisa nuestro valor y lo deja en 0pt. Si aquí se deja
  % distinto del footline normal, el pie de las láminas de sección queda
  % desalineado con el del resto.
  % Mismo \serifalab que el footline normal, y por la misma razón: solo en
  % «EMG | N/total», nunca en el nombre de la sección.
  \setbeamertemplate{footline}{%
    \hbox{%
      \begin{beamercolorbox}[wd=0.62\paperwidth, ht=2.6ex, dp=1.4ex, leftskip=1.4em]{}%
        {\usebeamerfont{footline}\color{white}\insertsectionhead}%
      \end{beamercolorbox}%
      \begin{beamercolorbox}[wd=0.38\paperwidth, ht=2.6ex, dp=1.4ex, right, rightskip=0.45em]{}%
        {\usebeamerfont{footline}\serifalab\color{verdeSombra}%
         EMG\hspace{0.9em}\textbar\hspace{0.9em}%
         \insertframenumber{}\,/\,\inserttotalframenumber}%
      \end{beamercolorbox}%
    }%
    \vskip2pt%
  }
  \begin{frame}[plain, noframenumbering]
    \vspace{1.4em}
    \begin{minipage}{0.86\textwidth}
      \hspace{1.2em}%
      \begin{minipage}{0.95\textwidth}
        \tableofcontents[currentsection, sectionstyle=show/shaded, hideallsubsections]
      \end{minipage}
    \end{minipage}
  \end{frame}
  }%
}

% ══════════════════════════════════════════════════════════════════════════════
% ENTORNOS PROPIOS DEL CURSO
% ══════════════════════════════════════════════════════════════════════════════

% Rótulo en versalitas falsas — tgheros sí tiene versalitas, pero el rótulo se
% ve mejor en mayúsculas con la letra pequeña.
\newcommand{\etiqueta}[1]{{\footnotesize\color{grisTenue}\MakeUppercase{#1}}}

% ── ¿Por qué nos importa esto? ────────────────────────────────────────────────
% Abre todas las sesiones: plantea la pregunta politológica antes de que
% aparezca una sola función.
\newcommand{\porquenosimporta}[2]{%
  \begin{frame}[plain, noframenumbering]
    \vfill
    \begin{center}
      \etiqueta{¿Por qué nos importa esto?}\\[1.3em]
      {\LARGE\color{verdeProfundo} #1}\\[1.5em]
      \begin{minipage}{0.80\textwidth}
        \centering\normalsize\color{grisTexto} #2
      \end{minipage}
    \end{center}
    \vfill
  \end{frame}%
}

% ── Lo que ya sabemos hacer ───────────────────────────────────────────────────
% Cierra todas las sesiones. Exactamente tres viñetas: es la recompensa.
\newcommand{\yasabemos}[3]{%
  \begin{frame}{Lo que ya sabemos hacer}
    \vspace{1.0em}
    \begin{itemize}\setlength{\itemsep}{1.4em}
      \item #1
      \item #2
      \item #3
    \end{itemize}
    \vfill
  \end{frame}%
}

% ── Adelanto de la sesión siguiente ───────────────────────────────────────────
\newcommand{\lasiguiente}[2]{%
  \begin{frame}[plain, noframenumbering]
    \vfill
    \begin{center}
      \etiqueta{La semana que viene}\\[1.3em]
      {\LARGE\color{verdeITAM} #1}\\[1.5em]
      \begin{minipage}{0.80\textwidth}
        \centering\normalsize\color{grisTexto} #2
      \end{minipage}
    \end{center}
    \vfill
  \end{frame}%
}

% ── Caja de error anticipado ──────────────────────────────────────────────────
% Los errores son parte del temario. Esta caja los presenta ANTES de que
% ocurran, con su traducción del inglés y su arreglo.
% "no shadow" y "sharp corners" son deliberados: nada de sombras ni degradados.
\newtcolorbox{errorbox}[1]{
  enhanced, breakable, no shadow, sharp corners,
  colback   = grisVelo,
  colframe  = guinda,
  boxrule   = 0pt, leftrule = 2.5pt,
  toprule = 0pt, bottomrule = 0pt, rightrule = 0pt,
  fontupper = \small\color{grisTexto},
  title     = {\color{guinda}\bfseries\small\ttfamily #1},
  colbacktitle = grisVelo,
  coltitle  = guinda,
  titlerule = 0pt
}

% ── Caja del salto politólogo ─────────────────────────────────────────────────
% El §7 pide que cada sesión contenga al menos un momento donde la intuición
% política se convierta en pregunta empírica, y que ese momento se nombre.
% Esta caja es el lugar donde se nombra.
\newtcolorbox{saltobox}{
  enhanced, breakable, no shadow, sharp corners,
  colback   = verdeVelo,
  colframe  = verdeITAM,
  boxrule   = 0pt, leftrule = 2.5pt,
  toprule = 0pt, bottomrule = 0pt, rightrule = 0pt,
  fontupper = \small\color{grisTexto},
  title     = {\color{verdeProfundo}\bfseries\small El salto},
  colbacktitle = verdeVelo,
  coltitle  = verdeProfundo,
  titlerule = 0pt
}

% ══════════════════════════════════════════════════════════════════════════════
% METADATOS FIJOS
% ══════════════════════════════════════════════════════════════════════════════
\author[E. Miranda]{Emiliano Miranda González}
\institute[ITAM]{%
  Instituto Tecnológico Autónomo de México\\
  Departamento Académico de Ciencia Política%
}

\hypersetup{
  colorlinks  = true,
  linkcolor   = verdeITAM,
  urlcolor    = verdeITAM,
  citecolor   = grisTenue,
  pdfauthor   = {Emiliano Miranda González},
  pdfsubject  = {Laboratorio de R para Ciencia Política — ITAM},
  pdfkeywords = {R, ciencia política, ITAM, métodos cuantitativos},
  unicode     = true
}
